\chapter{Neurology}

\section*{High-Yield Clinical Reference: Glasgow Coma Scale (GCS)}
\begin{center}
\small
\begin{tabular}{@{}lll@{}}
\toprule
\textbf{Category} & \textbf{Response Criteria} & \textbf{Score} \\
\midrule
\textbf{Eye Opening (E)} & Spontaneous & 4 \\
 & To sound & 3 \\
 & To pressure / pain & 2 \\
 & None & 1 \\
\midrule
\textbf{Verbal Response (V)} & Oriented & 5 \\
 & Confused & 4 \\
 & Words (inappropriate) & 3 \\
 & Sounds (incomprehensible) & 2 \\
 & None & 1 \\
\midrule
\textbf{Motor Response (M)} & Obeys commands & 6 \\
 & Localizes pain & 5 \\
 & Normal flexion (withdrawal) & 4 \\
 & Abnormal flexion (decorticate) & 3 \\
 & Extension (decerebrate) & 2 \\
 & None & 1 \\
\bottomrule
\end{tabular}
\end{center}
\noindent\textit{Total GCS Score = E + V + M (Range: 3 to 15). Severe brain injury: GCS $\le 8$ ("GCS 8, intubate").}
\bigskip

\medterm{Acoustic Neuroma}-- A benign, slow-growing schwannoma arising from the Schwann cells of the vestibular portion of the vestibulocochlear nerve (CN VIII), most commonly at the junction of the internal auditory canal and the cerebellopontine angle. Also known as vestibular schwannoma. Accounts for approximately 8 percent of intracranial tumors. Most are unilateral and sporadic; bilateral acoustic neuromas are pathognomonic of neurofibromatosis type 2 (NF2), an autosomal dominant condition caused by mutations in the NF2 gene on chromosome 22q12.2. Clinical presentation includes ipsilateral sensorineural hearing loss (typically high-frequency), tinnitus, vertigo or imbalance, and facial numbness or weakness from compression of the trigeminal or facial nerve. Diagnosis is confirmed by MRI with gadolinium, demonstrating an enhancing mass at the cerebellopontine angle with extension into the internal auditory canal. Treatment options include observation with serial imaging for small or slow-growing tumors, stereotactic radiosurgery (Gamma Knife), or microsurgical resection via a translabyrinthine, retrosigmoid, or middle cranial fossa approach.

\medterm{Age-Related Hearing Loss}-- Progressive, bilateral, symmetrical sensorineural hearing loss associated with aging, also known as presbycusis, affecting approximately one-third of adults between 65 and 74 years and nearly half of those over 75. The pathogenesis involves cumulative damage to cochlear structures including degeneration of outer hair cells (beginning at the basal turn affecting high-frequency sound transduction), loss of spiral ganglion neurons, stria vascularis atrophy (reducing endocochlear potential), and basilar membrane stiffness. Diagnosis is clinical, based on history of progressive bilateral high-frequency hearing loss with preserved speech discrimination in early stages. Audiometry demonstrates bilateral sensorineural hearing loss, typically with a downward-sloping configuration affecting higher frequencies more than lower frequencies. Management includes hearing aids as the mainstay of treatment, assistive listening devices, cochlear implantation for severe-to-profound loss, and avoidance of ototoxic medications.

\medterm{Agenesis of the Corpus Callosum}-- Complete or partial absence of the corpus callosum, the major commissural fiber bundle connecting the cerebral hemispheres. May be isolated or associated with other CNS anomalies (Dandy-Walker, Chiari, heterotopias) and genetic syndromes (Aicardi, Andermann, acrocallosal syndromes). Presents with a variable phenotype ranging from asymptomatic to developmental delay, intellectual disability, seizures, and impaired interhemispheric transfer. Diagnosis by prenatal or postnatal MRI (coronal and sagittal views show absent cingulate gyrus, colpocephaly, and radial orientation of medial sulci). No specific treatment; management addresses associated conditions.

\medterm{Agnosia}-- A higher-order neurological deficit characterized by the inability to recognize or process sensory stimuli (visual, auditory, tactile) despite intact primary sensory organs and pathways. Caused by focal lesions in specific cortical association areas (parietal, temporal, or occipital lobes). Types: Visual agnosia (inability to recognize objects visually), Prosopagnosia (inability to recognize familiar faces; fusiform gyrus lesion), Astereognosis (inability to identify objects by touch; contralateral parietal lobe lesion), and Anosognosia (unawareness or denial of one\'s own neurological deficit; right parietal lobe lesion).

\medterm{Agoraphobia}-- Anxiety disorder characterized by fear and avoidance of situations where escape might be difficult or help unavailable during a panic-like episode. Common feared situations include public transportation, open spaces, enclosed spaces, standing in line, being in crowds, and being outside the home alone. Often develops after panic attacks and becomes self-reinforcing through avoidance behavior. Severe cases result in housebound patients unable to leave home. Diagnosis: DSM-5 criteria. Treatment: cognitive behavioural therapy (CBT) as first-line psychological intervention, SSRIs (particularly sertraline, paroxetine, escitalopram) as first-line pharmacotherapy, and combination therapy for moderate-to-severe cases. Prognosis is variable, but with appropriate treatment most patients achieve significant improvement.

\medterm{Allodynia}-- A form of neuropathic pain where a normally non-painful stimulus evokes pain. It is a key feature of central and peripheral sensitization and can be categorized into tactile allodynia (pain from light touch or pressure), mechanical allodynia (pain from movement or stroking of the skin), and thermal allodynia (pain from normally non-painful temperature changes). Allodynia is commonly seen in postherpetic neuralgia, where the skin remains exquisitely sensitive to clothing or light touch. Diagnosis is clinical, supported by quantitative sensory testing and response to treatments. Management includes addressing the underlying cause, with pharmacological options including gabapentinoids (gabapentin, pregabalin), tricyclic antidepressants (amitriptyline, nortriptyline), serotonin-noradrenaline reuptake inhibitors (duloxetine, venlafaxine), lidocaine patches, capsaicin cream, and opioid analgesics as a last resort.

\medterm{Amyotrophic Lateral Sclerosis}-- Progressive neurodegenerative disorder affecting both upper and lower motor neurons, causing weakness, atrophy, fasciculations, and eventually respiratory failure. Most cases are sporadic (90 percent); 10 percent are familial. Average age of onset 55-60 years. Presents with asymmetric limb weakness (limb-onset) or bulbar symptoms (bulbar-onset: dysarthria, dysphagia). Upper motor neuron signs: spasticity, hyperreflexia, extensor plantar responses. Lower motor neuron signs: weakness, atrophy, fasciculations (visible muscle twitching), and cramps. Diagnosis is clinical (El Escorial criteria) with supportive electrodiagnostic studies (EMG showing active denervation and chronic reinnervation in multiple regions), neuroimaging to exclude mimics, and genetic testing for familial cases. Riluzole, a glutamate antagonist, prolongs survival by approximately 2-3 months by reducing excitotoxicity. Edaravone, a free radical scavenger, may slow functional decline in some patients. Multidisciplinary care including respiratory support (non-invasive ventilation for respiratory insufficiency), nutritional support (percutaneous endoscopic gastrostomy for dysphagia), speech therapy, and palliative care is essential for optimising quality of life. Median survival is 3-5 years from symptom onset, with respiratory failure as the most common cause of death.

\medterm{Anencephaly}-- A fatal neural tube defect where the cerebral hemispheres and cranial vault fail to develop due to failed closure of the rostral neuropore during the fourth week of gestation. The brainstem and spinal cord are present, and reflex movements may occur, but there is no cognitive function. Associated with polyhydramnios (due to impaired swallowing) and elevated maternal serum alpha-fetoprotein. Diagnosis by prenatal ultrasound. Most cases are stillborn or die within hours to days after birth. Prevention with folic acid supplementation.

\medterm{Anxiety Disorders}-- Group of psychiatric disorders characterized by excessive fear, worry, and anxiety that cause significant distress and impairment. Types include generalized anxiety disorder (persistent, excessive worry about multiple topics), panic disorder (recurrent unexpected panic attacks with palpitations, sweating, trembling), social anxiety disorder (fear of social situations), specific phobias (irrational fear of specific objects/situations), and agoraphobia (fear of situations where escape may be difficult or embarrassing. Treatment includes CBT as first-line psychological treatment and SSRIs (sertraline, escitalopram, paroxetine) as first-line pharmacotherapy. Prognosis is generally favourable with treatment, though relapse rates are significant.

\medterm{Aphasia}-- Impairment of language comprehension or production resulting from brain damage, most commonly from stroke affecting the left hemisphere (dominant) language areas. Types: Broca (expressive) aphasia from left inferior frontal gyrus damage—non-fluent speech, intact comprehension, frustration with communication; Wernicke (receptive) aphasia from left superior temporal gyrus damage—fluent but nonsensical speech, impaired comprehension, poor repetition; global aphasia from large left hemisphere stroke involving both anterior and posterior circulation territories —non-fluent speech, severely impaired comprehension, and poor repetition. Additional types include conduction aphasia (arcuate fasciculus lesion—fluent speech with impaired repetition but relatively intact comprehension), transcortical motor aphasia (disconnection of supplementary motor area—non-fluent speech with intact repetition), transcortical sensory aphasia (temporoparietal lesion—fluent speech with impaired comprehension but intact repetition), and anomic aphasia (angular gyrus lesion—word-finding difficulty with fluent speech). Diagnosis is through formal language assessment using the Boston Diagnostic Aphasia Examination or the Western Aphasia Battery. Treatment involves speech and language therapy, with compensatory strategies and communication aids for severe cases. Prognosis depends on the size and location of the lesion; Broca aphasia generally has a better recovery than Wernicke or global aphasia.

\medterm{Apraxia}-- Inability to perform learned, skilled motor movements despite intact motor function, sensation, coordination, and comprehension. Types: ideomotor apraxia (inability to perform movements on command, e.g., waving goodbye or using a hammer, due to parietal lobe lesions, typically left hemisphere), ideational apraxia (inability to sequence multi-step tasks, e.g., making a cup of tea, due to frontal or temporal lobe dysfunction), limb-kinetic apraxia (loss of fine motor dexterity, e.g., difficulty with finger tapping, due to premotor cortex lesions), buccofacial/oral apraxia (difficulty performing facial movements on command, e.g., blowing out a candle), and gait apraxia (difficulty initiating walking, common in normal pressure hydrocephalus). Apraxia is most commonly caused by stroke affecting the dominant (left) hemisphere, but can also result from Alzheimer disease, corticobasal degeneration, brain tumors, and traumatic brain injury. Diagnosis is clinical; there is no specific pharmacological treatment. Management: occupational therapy, task-specific training, and compensatory strategies.

\medterm{Ataxia}-- Lack of voluntary coordination of muscle movements, resulting from dysfunction of the cerebellum, posterior columns, or peripheral nerves. Types: cerebellar (gait and limb ataxia, dysarthria, nystagmus), sensory (impaired proprioception causing sensory ataxia, positive Romberg sign), and vestibular (vertigo, nystagmus, nausea). Cerebellar causes include stroke, tumors, alcohol, medications (phenytoin), degenerative diseases (spinocerebellar ataxias), and multiple sclerosis. Sensory causes include peripheral neuropathy, tabes dorsalis (neurosyphilis), and vitamin B12 deficiency causing posterior column dysfunction. Diagnosis is clinical, supported by neuroimaging (MRI brain for cerebellar causes, MRI spine for spinal causes), nerve conduction studies, and blood tests. Treatment is directed at the underlying cause, with symptomatic management including physical therapy, occupational therapy, and adaptive devices.

\medterm{Autism Spectrum Disorder}-- Neurodevelopmental disorder characterized by persistent deficits in social communication and social interaction across multiple contexts, restricted and repetitive patterns of behavior, interests, or activities. Symptoms present in early developmental period and cause clinically significant impairment. Severity levels: Level 1 (requiring support), Level 2 (requiring substantial support), Level 3 (requiring very substantial support). Associated features include intellectual disability (30-40 percent of cases), seizures, sleep disturbances, and gastrointestinal symptoms. Diagnosis is clinical using DSM-5 criteria, supported by multidisciplinary assessment including developmental history, cognitive testing, and speech and language evaluation. Management involves early intensive behavioural intervention (applied behaviour analysis), speech and language therapy, occupational therapy, social skills training, pharmacological treatment of associated conditions (SSRIs for anxiety, stimulants for inattention, antipsychotics for irritability and aggression), and educational support. Prognosis is highly variable, with early intervention being the most important predictor of functional outcome.

\medterm{Babinski Sign}-- The Babinski sign is a pathological reflex elicited by stroking the lateral aspect of the sole of the foot from the heel to the metatarsal heads with a blunt object. A normal response (flexor plantar response) is plantar flexion (curling downward) of all toes. An abnormal Babinski sign (extensor plantar response) is dorsiflexion (upgoing) of the great toe with fanning (abduction) of the other toes, indicating an upper motor neuron lesion affecting the corticospinal tract. The Babinski sign is a sign of corticospinal tract dysfunction rather than a specific disease. In infants, an extensor plantar response (upgoing toes) is normal until approximately 12-24 months of age due to incomplete myelination of the corticospinal tract, after which a flexor plantar response becomes established. A persistent Babinski sign beyond this age indicates upper motor neuron pathology and can be caused by stroke, multiple sclerosis, spinal cord injury, traumatic brain injury, cerebral palsy, amyotrophic lateral sclerosis, and other neurodegenerative conditions affecting the corticospinal tract.

\medterm{Balance}-- The ability to maintain the body's center of mass over its base of support, relying on input from the vestibular system, vision, and proprioception. It is essential for walking, standing, and performing daily activities without falling. Age-related decline, medications, and neurological disorders can impair balance, increasing fall risk.

\medterm{Balance Disorders and Vertigo}-- Balance disorders encompass conditions that cause dizziness, unsteadiness, or a sensation of motion, often arising from dysfunction in the inner ear, brain, or sensory pathways. Vertigo specifically describes the false sensation that the environment is spinning, commonly due to benign paroxysmal positional vertigo, vestibular neuritis, or M\'eni\`ere's disease. Diagnosis involves history, physical exam, and vestibular testing; treatment includes repositioning maneuvers, medications, and vestibular rehabilitation.

\medterm{Becker Muscular Dystrophy}-- An X-linked recessive neuromuscular disorder caused by mutations in the dystrophin gene, similar to Duchenne muscular dystrophy (DMD), but with in-frame mutations that allow production of some functional dystrophin protein, albeit reduced in quantity or abnormal in size. The result is a milder clinical phenotype compared to DMD. The incidence is approximately one in 18,000 to 30,000 male births. Clinical features include progressive proximal muscle weakness, typically beginning in the hip girdle and thighs (difficulty climbing stairs, rising from chairs, and walking), calf hypertrophy (pseudohypertrophy due to fatty infiltration, which is characteristic), waddling gait, Gowers manoeuvre (using hands to push off the floor and climb up the legs to stand), and reduced or absent deep tendon reflexes. Progression is slower than DMD, with loss of ambulation typically occurring in the fourth decade or later. Cardiac involvement (dilated cardiomyopathy) is common and is the leading cause of death, requiring regular cardiac surveillance. Respiratory function declines with disease progression. Serum creatine kinase levels are elevated (typically 5-10 times normal). Diagnosis is confirmed by genetic testing showing in-frame dystrophin gene mutations, with muscle biopsy reserved for cases where genetic testing is inconclusive. Management involves multidisciplinary care including physiotherapy to maintain flexibility, corticosteroids to slow progression (although less evidence than in DMD), cardiac monitoring with ACE inhibitors or beta-blockers, respiratory support, and orthopaedic management for contractures and scoliosis.

\medterm{Bell Palsy}-- Acute, idiopathic, unilateral facial nerve (CN VII) paralysis of lower motor neuron type, with an incidence of approximately 20-30 per 100,000 person-years. The cause is thought to be herpes simplex virus type 1 reactivation in the geniculate ganglion, leading to inflammation, demyelination, and edema of the facial nerve within the narrow confines of the facial canal (fallopian canal) of the temporal bone. Onset is acute and progressive over hours to days, with maximal weakness reached within 48 hours. Clinical features include acute-onset unilateral facial weakness that involves both the upper and lower face (forehead, eye closure, and mouth), inability to close the eye (lagophthalmos), which can lead to corneal exposure and ulceration, hyperacusis (due to stapedius muscle paralysis), loss of taste in the anterior two-thirds of the tongue (due to chorda tympani nerve involvement), and decreased tearing and salivation. Pain behind the ear is common and may precede the weakness. Diagnosis is clinical, requiring exclusion of other causes of facial paralysis such as stroke (upper motor neuron pattern sparing the forehead), Ramsay Hunt syndrome (herpes zoster oticus), Lyme disease, Guillain-Barre syndrome, otitis media, and parotid tumors. Treatment includes oral corticosteroids (prednisolone 60 mg daily for 7 days with rapid taper) started within 72 hours of symptom onset to improve outcomes, and eye protection with artificial tears, lubricating ointment, and taping the eye closed at night to prevent corneal damage. Antiviral agents (acyclovir or valacyclovir) are added if Ramsay Hunt syndrome is suspected. Prognosis is excellent, with approximately 70 percent of patients achieving near-complete recovery within 3-6 months, 15 percent developing some residual weakness, and 15 percent having moderate-to-severe permanent weakness.

\medterm{Bell's Palsy}-- See Bell Palsy.

\medterm{Brain Abscess}-- A focal, purulent infection within the brain parenchyma, most commonly occurring in the frontal or temporal lobes. Sources of infection include direct contiguous spread (otitis media and mastoiditis causing temporal lobe abscess; sinusitis causing frontal lobe abscess; dental infection), hematogenous dissemination from a distant source (bacterial endocarditis, lung abscess, congenital heart disease with right-to-left shunt), penetrating head trauma, or post-neurosurgical infection. Common causative organisms include Streptococcus milleri group (most common), anaerobic bacteria (Bacteroides, Fusobacterium), Staphylococcus aureus, and aerobic Gram-negative bacilli. Clinical presentation includes headache, fever, focal neurological deficits (hemiparesis, aphasia, visual field defects), seizures, and signs of increased intracranial pressure (nausea, vomiting, decreased level of consciousness). Diagnosis is by contrast-enhanced CT or MRI brain showing a ring-enhancing lesion with surrounding edema and central necrosis. Management includes empiric broad-spectrum antibiotics (third-generation cephalosporin plus metronidazole covering aerobic and anaerobic organisms), stereotactic aspiration or surgical excision for definitive diagnosis and source control, corticosteroids for vasogenic edema, and anticonvulsants for seizure prophylaxis. Mortality is approximately 10-20 percent, with prompt treatment improving outcomes.

\medterm{Brain Health}-- Brain health encompasses the ability to perform cognitive functions such as memory, learning, decision-making, and emotional regulation across the lifespan. Key factors include cardiovascular health, nutrition, sleep, social engagement, cognitive stimulation, and management of chronic conditions like hypertension and diabetes. Protecting brain health reduces the risk of dementia, stroke, and other neurological disorders.

\medterm{Brain Injuries}-- Brain injuries range from mild concussions to severe traumatic brain injury caused by falls, motor vehicle accidents, sports impacts, or penetrating trauma. Symptoms depend on injury location and severity, including headache, confusion, memory loss, seizures, and altered consciousness. Recovery varies widely and may involve rest, rehabilitation, cognitive therapy, and long-term supportive care.

\medterm{Brain Metastases}-- The most common intracranial neoplasms, occurring in approximately ten to forty percent of patients with systemic cancer, and are more common than primary brain tumors. They result from hematogenous spread of tumor cells to the brain, with the most common primary sources being lung cancer (the most common source overall, particularly non-small cell lung cancer and small cell lung cancer), breast cancer, melanoma (highest propensity for brain metastases), renal cell carcinoma, and colorectal cancer. Brain metastases typically occur at the gray-white matter junction, where the vascular supply changes from larger to smaller calibre vessels, and are most frequently located in the cerebral hemispheres (80 percent), cerebellum (15 percent), and brainstem (5 percent). Clinical presentation depends on the location and number of lesions but commonly includes headache, focal neurological deficits, seizures, cognitive impairment, and signs of increased intracranial pressure. Diagnosis is by contrast-enhanced MRI, which is more sensitive than CT for detecting small lesions. Management includes corticosteroids (dexamethasone) for vasogenic cerebral edema, whole-brain radiotherapy for multiple metastases, stereotactic radiosurgery for oligometastatic disease (typically 1-3 lesions), and surgical resection for single accessible lesions in patients with good performance status and controlled systemic disease.

\medterm{Brain Tumor}-- An abnormal growth of cells within the brain or its surrounding structures, classified as primary (originating from brain tissue or its coverings) or secondary/metastatic (spread from a cancer elsewhere in the body, which are more common, accounting for approximately 50 percent of all brain tumors). Primary brain tumors are classified by cell type and grade: meningiomas (most common primary brain tumor in adults, arising from the meninges, typically benign WHO grade 1), gliomas (arising from glial cells—including astrocytomas, oligodendrogliomas, glioblastomas (the most common and most malignant primary brain tumor, WHO grade 4), ependymomas, and mixed gliomas. Clinical presentation depends on the location, size, and rate of growth of the tumor but commonly includes progressive headache (often worse in the morning and worsened by Valsalva manoeuvre), seizures (focal or generalised), focal neurological deficits (hemiparesis, aphasia, visual field defects, sensory loss), cognitive and personality changes, and signs of increased intracranial pressure. Diagnosis is by neuroimaging: MRI with gadolinium is the imaging modality of choice, providing superior soft tissue resolution and demonstrating contrast enhancement, edema, mass effect, and relationship to adjacent structures. Definitive diagnosis requires histopathological examination of tissue obtained by stereotactic biopsy or surgical resection. Management is multidisciplinary and depends on tumor type, grade, and location, and includes surgical resection (gross total resection when safely achievable), radiotherapy, and chemotherapy (temozolomide for glioblastoma). Prognosis varies widely by tumor type and grade: benign meningiomas have excellent prognosis with surgical resection, while glioblastoma multiforme carries a poor prognosis with median survival of approximately 15 months despite aggressive multimodality treatment.

\medterm{Brain Ventricles}-- The four interconnected fluid-filled cavities within the brain that produce and contain cerebrospinal fluid (CSF). The lateral ventricles (first and second, paired, C-shaped, within each cerebral hemisphere, each with anterior/frontal horn, body, atrium, posterior/occipital horn, and temporal/inferior horn) communicate with the third ventricle via the interventricular foramina (foramina of Monro). The third ventricle (midline, within the diencephalon, between the thalami) communicates with the fourth ventricle via the cerebral aqueduct (aqueduct of Sylvius, within the midbrain). The fourth ventricle (between the brainstem and cerebellum) is continuous with the central canal of the spinal cord and drains CSF into the subarachnoid space via the midline foramen of Magendie and paired lateral foramina of Luschka. Hydrocephalus: accumulation of CSF causing ventricular enlargement, classified as communicating (impaired CSF absorption) or non-communicating/obstructive (blockage within the ventricular system). CSF is produced by the choroid plexus (modified ependymal cells) within the ventricles at a rate of ~500 mL/day., which encodes a copper-transporting ATPase in hepatocytes. The defective ATP7B protein impairs incorporation of copper into apoceruloplasmin and biliary excretion of copper, leading to progressive copper accumulation in the liver, brain, cornea, and other organs. Clinical manifestations typically present between ages five and thirty-five and include hepatic disease (ranging from asymptomatic elevation of liver enzymes to acute hepatitis, cirrhosis, and fulminant hepatic failure), neurological disease (dysarthria, ataxia, dystonia, parkinsonism, tremor—particularly a wing-beating tremor of the arms, chorea, and cognitive impairment with executive dysfunction), psychiatric symptoms (depression, personality changes, psychosis), and Kayser-Fleischer rings (golden-brown discolouration at the limbus of the cornea due to copper deposition in Descemet membrane, visible on slit-lamp examination and pathognomonic for Wilson disease when present). Diagnosis is based on a combination of clinical features, Kayser-Fleischer rings on slit-lamp examination, low serum caeruloplasmin (less than 0.2 g/L), elevated 24-hour urinary copper excretion (greater than 100 micrograms), elevated hepatic copper content on liver biopsy, and genetic testing for ATP7B mutations. Serum free copper (non-caeruloplasmin-bound copper) is elevated. Treatment involves lifelong copper chelation therapy with D-penicillamine (first-line, acts by chelating copper and promoting urinary excretion) or trientine, zinc acetate (which blocks intestinal absorption of copper and induces hepatic metallothionein, used as maintenance therapy or first-line for presymptomatic patients), and dietary restriction of copper-rich foods (shellfish, nuts, chocolate, mushrooms, organ meats). Liver transplantation is curative for severe hepatic failure or refractory neurological disease. Prognosis is excellent with early diagnosis and lifelong treatment; without treatment, Wilson disease is progressive and fatal. Patients who stop chelation therapy are at risk of irreversible neurological deterioration and hepatic decompensation.

\medterm{Brown Sequard Syndrome}-- Brown-Sequard syndrome (hemisection of the spinal cord) results from a lesion affecting one half of the spinal cord, producing a characteristic pattern of ipsilateral and contralateral neurological deficits. The ipsilateral side below the lesion demonstrates upper motor neuron paralysis (corticospinal tract), loss of vibration and proprioception (dorsal column), and loss of tactile sense. The contralateral side below the lesion demonstrates loss of pain and temperature sensation (spinothalamic tract). The syndrome is most commonly caused by penetrating trauma (knife or gunshot wounds), but may also result from spinal cord tumors, haematomas, disc herniation, demyelinating disease (multiple sclerosis), or infarction of the spinal cord (anterior spinal artery syndrome affecting the spinothalamic tract on the contralateral side). Diagnosis is clinical, confirmed by MRI of the spinal cord showing the lateralised lesion. Treatment is directed at the underlying cause and includes high-dose corticosteroids for inflammatory or demyelinating causes, surgical decompression for compressive lesions, and supportive care including pain management and rehabilitation. Prognosis depends on the etiology and severity of the cord injury, with some recovery possible over months.

\medterm{Carpal Tunnel Syndrome}-- The most common entrapment neuropathy, caused by compression of the median nerve as it passes through the carpal tunnel, a narrow fibro-osseous canal bounded by the carpal bones and the flexor retinaculum (transverse carpal ligament) at the wrist. The median nerve provides sensory innervation to the palmar aspect of the thumb, index finger, middle finger, and lateral half of the ring finger, and motor innervation to the thenar muscles (abductor pollicis brevis, opponens pollicis, and flexor pollicis brevis). The tunnel also contains nine flexor tendons (flexor digitorum profundus, flexor digitorum superficialis, and flexor pollicis longus) surrounded by synovial sheaths. Compression of the median nerve within the carpal tunnel results from any condition that reduces the volume of the tunnel or increases the volume of its contents: repetitive strain, obesity, pregnancy, diabetes, rheumatoid arthritis, hypothyroidism, acromegaly, amyloidosis, and wrist fractures. Clinical presentation includes numbness, tingling, and paraesthesiae in the median nerve distribution, typically occurring at night or with repetitive hand use. Patients may report waking with numb hands and shaking them for relief (flick sign). Motor weakness and thenar atrophy occur in advanced cases. Diagnosis is clinical, confirmed by nerve conduction studies showing prolonged distal motor and sensory latencies across the carpal tunnel. Treatment includes conservative measures (wrist splinting in neutral position at night, activity modification, NSAIDs, corticosteroid injection into the carpal tunnel) and surgical release (carpal tunnel release) for refractory or severe cases.

\medterm{Cauda Equina Syndrome}-- A neurosurgical emergency caused by compression of the cauda equina nerve roots below the level of the conus medullaris (typically at L1-L2 and below), resulting in bilateral lower extremity weakness, sensory loss, and bladder and bowel dysfunction. The most common cause is large lumbar disc herniation, but other causes include spinal tumors, epidural abscess, epidural hematoma, spinal stenosis, and trauma. Clinical features include severe low back pain, bilateral lower extremity weakness, bilateral sciatica, saddle anesthesia (loss of sensation in the perineal and perianal regions), loss of anal sphincter tone, urinary retention (painless, with overflow incontinence), and fecal incontinence. Urinary retention is an early and critical sign that is often missed; post-void residual volume greater than 100-200 mL is suggestive. Diagnosis is clinical, confirmed by urgent MRI of the lumbar spine showing compression of the cauda equina nerve roots. Management is a surgical emergency requiring urgent decompressive laminectomy and discectomy, ideally within 24-48 hours of symptom onset to maximise the chance of neurological recovery. Prognosis is better if surgery is performed early, particularly for urinary function, but long-term deficits may persist even with timely intervention.

\medterm{Central Nervous System (CNS)}-- The part of the nervous system comprising the brain and spinal cord, enclosed within the skull and vertebral column. The brain includes the cerebrum (consciousness, voluntary movement, sensation, language, cognition), cerebellum (coordination, balance, motor learning), and brainstem (vital autonomic centers, cranial nerve nuclei, ascending\/descending pathways). The spinal cord extends from the foramen magnum to L1--L2, transmitting motor and sensory information between brain and periphery via ascending (spinothalamic, dorsal column) and descending (corticospinal) tracts, and housing reflex arcs. The CNS is protected by the skull and vertebral column, meninges (dura mater, arachnoid mater, pia mater), and cerebrospinal fluid. The blood-brain barrier regulates molecular exchange between blood and neural tissue.

\medterm{Cerebellar Ataxia}-- A neurological sign characterized by impaired coordination of voluntary movements, resulting from dysfunction of the cerebellum or its connections. The cerebellum coordinates the timing, force, and accuracy of movements, and cerebellar dysfunction produces a characteristic constellation of findings. Gait ataxia presents with a wide-based, unsteady, staggering gait with difficulty performing tandem (heel-to-toe) walking. Limb ataxia includes dysmetria (overshooting or undershooting targets), dysdiadokokinesia (impaired rapid alternating movements), intention tremor (tremor that worsens as the target is approached), and rebound phenomenon. Speech is affected with scanning dysarthria (slow, irregular, slurred speech with variable volume and tempo). Ocular findings include nystagmus and impaired smooth pursuit. Cerebellar ataxia can be caused by stroke (ischemic or hemorrhagic), tumors, alcohol toxicity and nutritional deficiencies (thiamine), medications (phenytoin, carbamazepine, lithium, benzodiazepines), autoimmune conditions (anti-GAD, anti-glutamate receptor antibodies), paraneoplastic syndromes, neurodegenerative diseases (spinocerebellar ataxias, multiple system atrophy), multiple sclerosis, infectious and post-infectious causes, and hypothyroidism. Diagnosis is clinical with MRI brain to identify structural, inflammatory, or degenerative causes. Management focuses on treating the underlying cause and providing symptomatic relief with physical and occupational therapy, and assistive devices for gait instability.

\medterm{Cerebral Aneurysm}-- An abnormal focal dilation of a cerebral artery, most commonly at bifurcations of the Circle of Willis. Saccular (berry) aneurysms account for 90\% of intracranial aneurysms. Risk factors: smoking, hypertension, female sex, polycystic kidney disease, Ehlers-Danlos, family history. Most are asymptomatic (incidental finding). Rupture causes subarachnoid hemorrhage: thunderclap headache, nuchal rigidity, photophobia, and altered mental status. Diagnosis: CTA or MRA for detection; DSA is gold standard. Treatment: microsurgical clipping or endovascular coiling for ruptured and high-risk unruptured aneurysms.

\medterm{Cerebral Cortex}-- The outermost layer of the cerebrum, composed of gray matter (neuronal cell bodies) folded into gyri (ridges) and sulci (grooves), greatly increasing surface area. Divided into four lobes: frontal (motor cortex, prefrontal cortex--executive function, personality, Broca area--speech production), parietal (somatosensory cortex, spatial orientation, integration of sensory modalities), temporal (auditory cortex, Wernicke area--language comprehension, memory formation via hippocampus), and occipital (primary visual cortex, visual processing). The cortex is organized into six horizontal layers (molecular, external granular, external pyramidal, internal granular, internal pyramidal, multiform) with distinct cell types and connectivity. Cortical functions are topographically organized (homunculus in motor and sensory cortices). Damage causes focal deficits corresponding to the affected region.

\medterm{Cerebrospinal Fluid (CSF)}-- A clear, colorless fluid that fills the ventricles of the brain, the central canal of the spinal cord, and the subarachnoid space around the brain and spinal cord. Produced primarily by the choroid plexus of the lateral, third, and fourth ventricles (~500 mL\/day, total volume ~150 mL in adults). Functions: mechanical cushioning of the brain (buoyancy reduces effective brain weight by ~97\%), removal of metabolic waste, transport of hormones and nutrients, and maintenance of stable intracranial pressure. Normal composition: protein 15--45 mg\/dL, glucose 50--80 mg\/dL (two-thirds of serum glucose), WBC 0--5\/mcL (mononuclear). CSF analysis by lumbar puncture is essential for diagnosing meningitis, subarachnoid hemorrhage, multiple sclerosis, and other neurologic conditions.

\medterm{Cerebrovascular Accident (Stroke)}-- Stroke is a neurological deficit caused by focal cerebral ischemia (ischemic stroke, 87\%) or hemorrhage (hemorrhagic stroke, 13\%). Ischemic stroke results from thrombotic or embolic occlusion of a cerebral artery. Risk factors: hypertension, atrial fibrillation, diabetes, smoking. Management: acute thrombolysis (tPA within 4.5 hours), mechanical thrombectomy for large vessel occlusion, secondary prevention with antiplatelets and risk factor control.

\medterm{Cervical Disc Herniation}-- A common cause of neck pain and radiculopathy, occurring when the nucleus pulposus of an intervertebral disc in the cervical spine herniates through the annulus fibrosus, compressing adjacent nerve roots or the spinal cord. The most commonly affected levels are C5-C6 and C6-C7, accounting for approximately ninety-five percent of cervical disc herniations. Cervical radiculopathy presents with dermatomal pain radiating from the neck into the shoulder, arm, or hand, along with sensory changes (numbness, tingling) and motor weakness in the corresponding myotome. C5-C6 disc herniation compresses the C6 nerve root, causing pain radiating to the lateral arm, thumb, and index finger, weakness of biceps, brachioradialis, and wrist extensors, and decreased biceps reflex. C6-C7 disc herniation compresses the C7 nerve root, causing pain radiating to the middle finger, weakness of triceps and finger extensors, and decreased triceps reflex. Diagnosis is by MRI of the cervical spine, which demonstrates the herniated disc and its relationship to the nerve root or spinal cord. Electromyography and nerve conduction studies may confirm the radiculopathy and exclude other causes. Management includes conservative treatment (analgesics, NSAIDs, muscle relaxants, physical therapy, cervical immobilisation) for most cases, with surgical intervention (anterior cervical discectomy and fusion, cervical disc replacement) indicated for progressive motor deficit, severe refractory pain, or spinal cord compression with myelopathy.

\medterm{Cervical Radiculopathy}-- A condition caused by compression, irritation, or inflammation of a spinal nerve root in the cervical spine, resulting in pain, numbness, tingling, and weakness in the distribution of the affected nerve root. The most commonly affected nerve roots are C6 and C7, accounting for approximately seventy to eighty percent of cervical radiculopathies. C5 radiculopathy causes shoulder pain, weakness of deltoid and biceps, and decreased biceps reflex. C6 radiculopathy causes pain radiating from the neck into the lateral arm, thumb, and index finger, weakness of biceps and brachioradialis, and decreased biceps reflex. C7 radiculopathy causes pain radiating into the middle finger, weakness of triceps and wrist and finger extensors, and decreased triceps reflex. C8 radiculopathy causes pain radiating into the medial arm, ring and little fingers, weakness of hand intrinsics, and decreased finger jerk. Diagnosis is by MRI of the cervical spine, with EMG/NCS confirming the level of nerve root involvement and excluding peripheral nerve entrapment. Management is initially conservative with analgesics, NSAIDs, physical therapy, and cervical epidural steroid injections. Surgical options include anterior cervical discectomy and fusion (ACDF) or cervical disc replacement for persistent or progressive symptoms despite conservative therapy, or for cases with myelopathy.

\medterm{Chiari Malformation}-- Chiari malformations are structural abnormalities of the posterior fossa and cerebellum characterized by downward displacement of the cerebellar tonsils through the foramen magnum. Type I is the most common variant, defined as cerebellar tonsillar ectopia of five millimeters or more below the foramen magnum. It may be asymptomatic or present with occipital headache (worsened by Valsalva maneuver, coughing, or straining), neck pain, lower cranial nerve palsies, cerebellar ataxia, myelopathy, and syringomyelia (a fluid-filled cavity within the spinal cord). Type II (Arnold-Chiari malformation) is a more severe variant associated with myelomeningocoele and hydrocephalus. Diagnosis is made by MRI of the brain and cervical spine in the sagittal plane, which demonstrates tonsillar herniation. Surgical decompression via suboccipital craniectomy and cervical laminectomy is indicated for symptomatic patients with brainstem compression, syringomyelia, or progressive neurological deficits.

\medterm{Chronic Fatigue Syndrome}-- A complex, debilitating disorder characterized by profound fatigue lasting >6 months that is not relieved by rest and is worsened by physical or mental exertion (post-exertional malaise). Associated symptoms: unrefreshing sleep, cognitive impairment (brain fog), orthostatic intolerance, joint pain, headaches, and sore throat. Also known as myalgic encephalomyelitis (ME/CFS). Diagnosis: clinical (CDC or Canadian Consensus criteria); exclusion of other causes. No FDA-approved treatment; management focuses on symptom relief, graded activity pacing, cognitive behavioral therapy for coping, and treating comorbid conditions (sleep disorders, orthostatic intolerance).

\medterm{Chronic Inflammatory Demyelinating Polyneuropathy}-- An acquired, immune-mediated peripheral neuropathy with a relapsing-remitting or progressive course. Presents with symmetric proximal and distal weakness, sensory loss, areflexia, and elevated CSF protein. Distinguishing from Guillain-Barre: progressive course >8 weeks (vs. rapid onset in GBS). Diagnosis: EMG/NCS shows demyelinating features (conduction block, temporal dispersion). Treatment: corticosteroids, IVIG, or plasmapheresis. Immunosuppressants (azathioprine, mycophenolate, rituximab) for refractory cases.

\medterm{CIDP}-- Chronic inflammatory demyelinating polyneuropathy (CIDP) is a chronic, immune-mediated, peripheral neuropathy characterized by progressive or relapsing symmetrical weakness and sensory loss in all four limbs, with areflexia or hyporeflexia, lasting for more than eight weeks. It is the chronic counterpart of GBS and is considered to be caused by an antibody-mediated attack on peripheral nerve myelin. The pathophysiology involves both humoral and cellular immune mechanisms targeting myelin and Schwann cells. Clinical presentation includes symmetrical proximal and distal weakness in all four limbs, with early loss of deep tendon reflexes, sensory loss (impaired vibration and proprioception), and symptoms developing progressively over weeks to months. Diagnosis is made using the European Federation of Neurological Societies (EFNS) criteria, supported by nerve conduction studies showing demyelinating features (prolonged distal latencies, reduced conduction velocities, conduction block, temporal dispersion, prolonged F-wave latencies), CSF analysis showing elevated protein with normal cell count (cytoalbuminological dissociation), and MRI of the brachial or lumbosacral plexus showing nerve root thickening and gadolinium enhancement. Management includes corticosteroids (prednisolone), intravenous immunoglobulin (IVIG), and plasma exchange as first-line immunomodulatory therapies. Immunosuppressive agents (azathioprine, mycophenolate mofetil, cyclophosphamide, rituximab) are used for refractory or relapsing cases. Prognosis is variable; most patients require long-term immunomodulatory therapy to maintain function and prevent disability progression.

\medterm{Cluster Headache}-- A primary headache disorder characterized by severe, unilateral, strictly orbital or periorbital pain attacks lasting fifteen to one hundred eighty minutes, occurring with one of three autonomic symptoms on the same side as the pain: conjunctival injection, lacrimation, nasal congestion or rhinorrhea, eyelid edema, forehead or facial sweating, miosis, or ptosis. Attacks occur in clusters or cycles, typically one to eight attacks per day during a cluster period lasting weeks to months followed by remission periods lasting months to years. Cluster headache is the most severe primary headache disorder, often described as the most painful condition known to medicine, and is nicknamed suicide headache due to its intensity. The headache is accompanied by restlessness and agitation—patients typically pace the floor, rock back and forth, or bang their head during attacks, in contrast to migraine patients who prefer to lie still in a dark room. Diagnosis is clinical using ICHD-3 criteria. Acute treatment includes high-flow oxygen (100\% at 12-15 L/min via non-rebreather mask) and subcutaneous sumatriptan or intranasal zolmitriptan, which cause vasoconstriction of dilated orbital vessels. Preventive treatment includes verapamil (first-line), lithium, topiramate, galcanezumab (a CGRP monoclonal antibody specifically approved for cluster headache), and greater occipital nerve blockade. Transitional therapy with oral or intravenous corticosteroids is used at the onset of a cluster period to break the cycle while long-term preventatives take effect.

\medterm{Cognitive Fitness}-- The active maintenance of mental sharpness and brain health through stimulating activities, physical exercise, social engagement, and proper nutrition. Regular challenges such as learning new skills, puzzles, reading, and social interaction support neuroplasticity and cognitive reserve. Maintaining cognitive fitness throughout life may delay age-related cognitive decline and reduce dementia risk.

\medterm{Cognitive Health and Memory}-- Cognitive health encompasses the ability to think clearly, learn, remember, and make decisions, with memory being a core component that involves encoding, storing, and retrieving information. Age-related changes in processing speed and recall are normal, but persistent decline may signal mild cognitive impairment or dementia. Protective factors include cardiovascular health, mental stimulation, physical activity, and social engagement.

\medterm{Coma}-- A state of profound unconsciousness in which the patient cannot be aroused and shows no awareness of self or environment, with eyes closed and no sleep-wake cycles. Results from widespread bilateral cerebral hemisphere dysfunction or brainstem reticular formation damage. Causes: structural (traumatic brain injury, stroke, intracranial hemorrhage, brain tumor), metabolic (hypoglycemia, hepatic encephalopathy, uremia, hyper\/hyponatremia, hypothyroidism), toxic (drug overdose--opioids, benzodiazepines, alcohol, barbiturates), infectious (meningitis, encephalitis), and hypoxic-ischemic (cardiac arrest). Assessment: Glasgow Coma Scale (eye, motor, verbal responses; range 3--15). Management: stabilize ABCs, treat reversible causes (naloxone, flumazenil, glucose, thiamine), neuroimaging, EEG to rule out nonconvulsive status epilepticus. Prognosis depends on cause, duration, and neurologic examination findings.

\medterm{Communicating Hydrocephalus}-- A type of hydrocephalus in which there is free flow of cerebrospinal fluid (CSF) between the ventricles, but impaired absorption of CSF at the arachnoid granulations or increased CSF production. The ventricles are uniformly dilated, including the fourth ventricle. The most common cause is impaired CSF absorption due to prior subarachnoid hemorrhage, meningitis, or intracranial surgery, which causes scarring and fibrosis of the arachnoid granulations. Other causes include idiopathic intracranial hypertension, cerebral venous sinus thrombosis, and overproduction of CSF by a choroid plexus papilloma. Clinical presentation includes headache, nausea and vomiting, gait disturbance, cognitive slowing, and papilloedema from increased intracranial pressure. In contrast to obstructive (non-communicating) hydrocephalus, the ventricles are uniformly dilated and there is no obstruction to CSF flow within the ventricular system. Diagnosis is by neuroimaging (CT or MRI) showing ventricular dilation with no obstructive lesion, supported by CSF drainage or infusion studies. Management includes CSF diversion procedures such as ventriculoperitoneal (VP) shunting or endoscopic third ventriculostomy (ETV) in selected cases where an obstructive component is present. The underlying cause must also be addressed (e.g., treatment of meningitis, management of subarachnoid hemorrhage complications).

\medterm{Complex Regional Pain Syndrome}-- A chronic pain condition characterized by regional pain and sensory changes disproportionate to the inciting event. Type I (reflex sympathetic dystrophy): no identifiable nerve injury. Type II (causalgia): follows a major nerve injury. Presents with burning pain, allodynia, hyperalgesia, edema, vasomotor changes (temperature and color changes), sudomotor changes (sweating), and trophic changes (skin, hair, nail changes). Diagnosis: Budapest clinical criteria. Treatment: early physical therapy and occupational therapy; NSAIDs, gabapentinoids, tricyclic antidepressants, bisphosphonates; sympathetic nerve blocks for severe pain.

\medterm{Concussion}-- A form of mild traumatic brain injury caused by a biomechanical force (direct blow to the head, face, neck, or elsewhere on the body with impulsive force transmitted to the head), resulting in a transient disturbance of brain function without structural damage visible on conventional neuroimaging. Concussion is a functional rather than structural injury, involving neurometabolic dysfunction: neuronal membrane disruption, ionic shifts (potassium efflux, calcium influx), and increased glucose metabolism (initially hyperglycolysis followed by metabolic depression), lactate accumulation, and impaired mitochondrial function. The result is a state of metabolic vulnerability in which the brain is less able to tolerate a second impact, which is the basis for the principle of avoiding repeat head injury during the recovery period. Clinical features include headache, dizziness, nausea, balance problems, blurred vision, sensitivity to light and noise, fatigue, sleep disturbance, difficulty concentrating, memory impairment, and emotional lability (irritability, anxiety, depression). Loss of consciousness occurs in fewer than 10 percent of concussions; post-traumatic amnesia is common. Diagnosis is clinical, based on history of head trauma and symptom assessment using standardised tools such as the Sport Concussion Assessment Tool (SCAT-5) or the Standardised Assessment of Concussion (SAC). Neuroimaging (CT or MRI) is normal. Management includes physical and cognitive rest during the acute phase (24-48 hours), followed by gradual stepwise return to activity guided by symptom tolerance. Most patients recover fully within 7-14 days in adults and within 2-4 weeks in children and adolescents. Persistent symptoms beyond 4 weeks are termed post-concussion syndrome and require multidisciplinary management.

\medterm{Conus Medullaris Syndrome}-- A clinical condition resulting from damage to the conus medullaris, the tapered, terminal end of the spinal cord located at approximately the L1-L2 vertebral level. The conus medullaris contains the sacral spinal cord segments (S2-S5) and the lowest lumbar segments, which innervate the bladder, bowel, sexual function, and the lower extremities. Clinical features include acute or subacute bilateral lower extremity weakness (usually symmetric), saddle anesthesia (loss of sensation in the perineal, perianal, and inner thigh regions), and early bladder and bowel dysfunction (urinary retention, overflow incontinence, and fecal incontinence). The presence of a bulbocavernosus reflex and anal wink reflex helps distinguish conus medullaris syndrome from cauda equina syndrome, and the early onset of bladder and bowel dysfunction (often before significant lower extremity weakness) is a distinguishing feature. Causes include compression from tumors (metastases, ependymomas), trauma, lumbar disc herniation, spinal cord infarction from anterior spinal artery occlusion, and inflammatory conditions such as transverse myelitis. Diagnosis is by urgent MRI of the thoracolumbar spine. Management requires prompt treatment of the underlying cause, including surgical decompression for compressive lesions. Prognosis is often poor for neurological recovery, particularly for bladder and bowel function.

\medterm{Corticobasal Degeneration}-- A rare, progressive neurodegenerative disorder characterized by asymmetric cortical and basal ganglia dysfunction. Presents with asymmetric parkinsonism (rigidity, bradykinesia), apraxia (loss of learned motor skills), alien limb phenomenon (involuntary, purposeless movements), cortical sensory loss, myoclonus, and dystonia. Later: cognitive decline, behavioral changes, and aphasia. Diagnosis: clinical; MRI may show asymmetric frontoparietal atrophy. No disease-modifying treatment; symptomatic therapy with levodopa (poor response), botulinum toxin for dystonia, and supportive care.

\medterm{Creutzfeldt Jakob Disease}-- Creutzfeldt-Jakob disease (CJD) is a rare, fatal, transmissible spongiform encephalopathy caused by the accumulation of misfolded prion protein (PrPSc) in the brain, leading to rapid, progressive dementia and myoclonus. Sporadic CJD is the most common form, accounting for approximately eighty-five percent of cases, with an incidence of approximately one to two per million per year, typically affecting patients in their sixties. The clinical presentation includes rapidly progressive cognitive decline, myoclonus (brief, involuntary muscle jerks that are pathognomonic when elicited by startle), ataxia, visual disturbances (cortical blindness, visual hallucinations), and pyramidal and extrapyramidal signs. In the late stages, patients become akinetic and mute. The classic EEG finding is generalised periodic sharp-wave complexes (PSWCs) at 1-2 Hz, which are present in approximately two-thirds of cases. MRI brain may show hyperintensity in the basal ganglia (particularly the caudate and putamen) and cortical ribboning on diffusion-weighted imaging. CSF analysis may detect 14-3-3 protein and elevated tau protein, and real-time quaking-induced conversion (RT-QuIC) testing of CSF has high sensitivity and specificity for detecting PrPSc. The disease is uniformly fatal, usually within 6-12 months of symptom onset. Variant CJD (vCJD), linked to bovine spongiform encephalopathy, affects younger patients and has a longer disease course. Management is supportive, with no effective disease-modifying treatment available.

\medterm{CVA}-- See Cerebrovascular Accident (Stroke).

\medterm{Dandy-Walker Malformation}-- A congenital hindbrain malformation characterized by hypoplasia or aplasia of the cerebellar vermis, cystic dilation of the fourth ventricle, and posterior fossa enlargement with upward displacement of the tentorium, transverse sinuses, and torcular herophili. Presents with hydrocephalus, developmental delay, ataxia, and cranial nerve palsies. Associated with agenesis of the corpus callosum and other CNS anomalies. Diagnosis by prenatal or postnatal MRI. Treatment: VP shunt or endoscopic third ventriculostomy for hydrocephalus. Prognosis depends on associated anomalies.

\medterm{Diabetic Peripheral Neuropathy}-- The most common complication of diabetes mellitus, affecting approximately fifty percent of patients with long-standing disease. It is a length-dependent, sensorimotor polyneuropathy that affects the distal lower extremities first, in a glove-and-stocking distribution, and progresses proximally over time. The pathophysiology involves multiple mechanisms including metabolic injury from hyperglycemia (polyol pathway, advanced glycation end products, oxidative stress, mitochondrial dysfunction, activation of the poly(ADP-ribose) polymerase pathway, and impaired axonal transport), microvascular disease (endoneurial capillary damage causing endoneurial hypoxia), and inflammatory mechanisms. Clinical presentation begins with numbness, tingling, and paraesthesiae in the toes and feet, progressing proximally over years. Pain is a common and disabling feature, often described as burning, shooting, or stabbing. Motor weakness (foot drop, intrinsic hand muscle weakness) develops in advanced stages. Physical examination reveals reduced vibration sense (the earliest finding), impaired pinprick and temperature sensation in a stocking distribution, loss of ankle reflexes, and eventually proprioceptive loss causing sensory ataxia. Autonomic involvement can cause orthostatic hypotension, gastroparesis, erectile dysfunction, and abnormal sweating. Diagnosis is clinical, supported by nerve conduction studies showing a length-dependent axonal sensorimotor polyneuropathy. Management involves optimising glycaemic control (the only intervention that slows progression), foot care to prevent ulceration, and symptomatic treatment of neuropathic pain with gabapentinoids, tricyclic antidepressants, or SNRIs (duloxetine).

\medterm{Duchenne Muscular Dystrophy}-- An X-linked recessive neuromuscular disorder caused by mutations in the dystrophin gene (DMD gene) on chromosome Xp21.2, resulting in absent or severely reduced dystrophin protein, which is essential for maintaining the structural integrity of the sarcolemma during muscle contraction. Most mutations are frameshift mutations (deletions, duplications, or point mutations) that introduce premature stop codons, leading to truncated, non-functional dystrophin. The incidence is approximately one in 3,500 male births, making it the most common muscular dystrophy in children. Onset of symptoms typically occurs before age five, with delayed motor milestones (late walking), difficulty running and climbing stairs, toe walking, and a waddling gait. Gowers manoeuvre (using the hands to push against the thighs to stand from a sitting position) is a classic finding. Calf pseudohypertrophy (enlargement due to fatty infiltration) is characteristic. Progression is rapid: most boys lose ambulation by age 12, develop respiratory insufficiency requiring non-invasive ventilation in adolescence, and die from respiratory or cardiac failure in their twenties or thirties. Cardiac involvement (dilated cardiomyopathy) is universal and requires regular surveillance with echocardiography and ECG. Diagnosis is confirmed by genetic testing showing out-of-frame dystrophin gene deletions, duplications, or point mutations, or by muscle biopsy showing absent dystrophin staining. Management includes corticosteroids (prednisolone or deflazacort) to prolong ambulation and reduce the need for scoliosis surgery, angiotensin-converting enzyme inhibitors for cardiomyopathy management, respiratory support, physiotherapy, and orthopaedic management for contractures and scoliosis. Gene therapy (delandistrogene moxeparvovec) has recently been approved for ambulatory patients with specific DMD mutations, and exon skipping agents (eteplirsen, golodirsen, viltolarsen) are available for specific mutation subtypes.

\medterm{DWI FLAIR Mismatch}-- Diffusion-weighted imaging (DWI) and fluid-attenuated inversion recovery (FLAIR) MRI sequences are used together to estimate the age of an ischemic stroke lesion and guide treatment decisions. DWI demonstrates restricted diffusion in acute ischemic tissue within minutes of stroke onset, appearing as a hyperintense (bright) lesion with corresponding hypointensity on the apparent diffusion coefficient (ADC) map. FLAIR demonstrates increased signal from vasogenic edema, which takes approximately four to six hours to become visible on FLAIR after stroke onset. A DWI-positive, FLAIR-negative lesion, indicating that the stroke is within the window where the patient may benefit from thrombolysis even if the onset time is unknown (e.g., wake-up stroke). If both DWI and FLAIR are positive, the lesion is estimated to be greater than approximately 4.5 hours old and thrombolysis is relatively contraindicated. This mismatch paradigm has been validated in the WAKE-UP trial, which demonstrated that patients with DWI-FLAIR mismatch who received alteplase had better functional outcomes at 90 days compared to placebo, without increased mortality.

\medterm{Dysarthria}-- A motor speech disorder resulting from neurological impairment of the muscular components of speech production, including the lips, tongue, palate, vocal cords, and respiratory muscles. It is caused by lesions affecting the central or peripheral nervous system at various levels: upper motor neuron (stroke, traumatic brain injury, multiple sclerosis), lower motor neuron (bulbar palsy, myasthenia gravis, Guillain-Barre syndrome), extrapyramidal (Parkinson disease, Huntington disease), cerebellar (ataxic dysarthria from cerebellar lesions), and neuromuscular (myasthenia gravis, muscular dystrophies, motor neuron disease). The type of dysarthria can help localise the neurological lesion and guide diagnosis. Flaccid dysarthria (lower motor neuron) presents with hypernasal, breathy, imprecise speech with audible inspiration. Spastic dysarthria (upper motor neuron) presents with strained, strangled, slow speech with reduced range of movement. Ataxic dysarthria (cerebellar) presents with scanning, irregular, and explosive speech. Hypokinetic dysarthria (extrapyramidal, e.g., Parkinson disease) presents with quiet, rapid, monotonous, mumbling speech with festination. Hyperkinetic dysarthria (e.g., Huntington disease) presents with variable speech rate, involuntary vocalisations, and variable volume. Diagnosis is by speech-language pathology assessment and identification of the underlying neurological cause. Management includes speech and language therapy, treatment of the underlying neurological condition, and augmentative and alternative communication devices for severe cases.

\medterm{Dysesthesia}-- An unpleasant, abnormal sensation produced by a normally non-painful stimulus, or an unpleasant, exaggerated response to a normal stimulus. It is distinguished from paresthesia (which is simply an abnormal sensation that may or may not be unpleasant) by its consistently unpleasant quality. Dysesthesia is a feature of neuropathic pain and is caused by dysfunction of the somatosensory nervous system, including peripheral nerve damage, spinal cord lesions, or central brain lesions. Peripheral causes include compression neuropathies, radiculopathies, and polyneuropathies. The distinction between dysesthesia and other sensory symptoms is important for diagnosis: dysesthesia is consistently unpleasant (a key distinguishing feature from paresthesia), allodynia is pain from a normally non-painful stimulus, hyperalgesia is an exaggerated response to a normally painful stimulus, and hyperpathia is a delayed, explosive, and unpleasant response to a stimulus, often with radiation beyond the stimulus site. Diagnosis is clinical, with neurological examination and nerve conduction studies to identify the underlying somatosensory pathway lesion. Management includes treating the underlying cause and symptomatic treatment of the dysesthetic pain with gabapentinoids, tricyclic antidepressants, or SNRIs.

\medterm{Dysphasia}-- A language disorder caused by damage to the language-dominant hemisphere of the brain (typically left), involving difficulty with expression, comprehension, reading, or writing. Dysphasia (often used interchangeably with aphasia): partial loss of language function. Types: expressive dysphasia (Broca area--non-fluent, effortful speech with preserved comprehension, agrammatism), receptive dysphasia (Wernicke area--fluent but meaningless speech, poor comprehension, word substitutions (paraphasias)), global dysphasia (both hemispheres extensively affected, all language modalities impaired), conduction dysphasia (arcuate fasciculus lesion--good comprehension and fluency but poor repetition), and anomic dysphasia (word-finding difficulty with otherwise fluent speech). Most common cause: ischemic stroke. Management: speech and language therapy.

\medterm{Dystonia}-- A movement disorder characterized by sustained or intermittent muscle contractions causing abnormal, often repetitive, movements and postures. It can be classified by the body region affected (focal, segmental, multifocal, or generalized), by age of onset (childhood, adolescent, or adult), and by etiology (primary/genetic or secondary/symptomatic). Focal dystonia includes cervical dystonia (torticollis, the most common adult-onset focal dystonia), blepharospasm (involuntary closure of the eyelids (blepharospasm), oromandibular dystonia (jaw opening, closing, or deviation), laryngeal dystonia (spasmodic dysphonia), and task-specific dystonias such as writer's cramp and musician's dystonia. Secondary dystonia can be caused by drugs (tardive dystonia from antipsychotics), stroke, traumatic brain injury, cerebral palsy, Wilson disease, and neurodegenerative disorders. Diagnosis is clinical, supported by genetic testing for primary dystonia (DYT genes including TOR1A, THAP1, GCH1, GNAL) and neuroimaging to exclude structural causes. Management includes botulinum toxin injections for focal dystonia (the treatment of choice), oral medications (anticholinergics, benzodiazepines, baclofen, tetrabenazine), and deep brain stimulation of the globus pallidus internus for severe, medically refractory generalised or segmental dystonia.

\medterm{Echolalia}-- The involuntary, automatic repetition of words or phrases spoken by another person, a speech phenomenon observed in various neurologic and psychiatric conditions. Immediate echolalia involves repeating sounds immediately after hearing them; delayed echolalia involves repeating phrases heard hours, days, or weeks earlier. Most commonly associated with autism spectrum disorder (where it may serve communicative or self-regulatory functions), Tourette syndrome, and transcortical sensory aphasia. Also occurs in Alzheimer disease, frontotemporal dementia, catatonia, and schizophrenia. Echolalia can be distinguished from palilalia (repetition of one's own words) and coprolalia (involuntary utterance of obscenities).

\medterm{Electroencephalography}-- A non-invasive neurophysiological test that records the electrical activity of the brain using electrodes placed on the scalp according to the international 10-20 system. The EEG measures the summated postsynaptic potentials of cortical pyramidal neurons and displays them as waveforms characterized by frequency (delta 0.5-4 Hz, theta 4-8 Hz, alpha 8-13 Hz, beta 13-30 Hz, gamma greater than 30 Hz), amplitude, and morphology. The normal waking EEG shows a posterior dominant alpha rhythm (8-13 Hz) that attenuates with eye opening. EEG is essential for the diagnosis of epilepsy, classification of seizure types and syndromes, and monitoring of status epilepticus. Abnormal findings include epileptiform discharges (spikes, sharp waves, spike-wave complexes) indicating an epileptic focus or generalised epilepsy. EEG is also used in the evaluation of encephalopathy (generalised slowing), brain death (electrocerebral silence), and sleep disorders. Prolonged video-EEG monitoring is used for presurgical evaluation of drug-resistant epilepsy, capturing ictal and interictal EEG patterns to localise the seizure onset zone. Standard EEG records from 21 scalp electrodes at the standard 10-20 positions, with additional electrodes (T1/T2, zygomatic, sphenoidal) used for better localisation of temporal lobe foci. Limitations include low spatial resolution, sensitivity to movement artifact, and inability to detect deep or mesial structures.

\medterm{Encephalitis}-- Inflammation of the brain parenchyma, most commonly caused by viral infections (herpes simplex virus, arboviruses, enteroviruses). Presents with fever, headache, altered mental status, seizures, and focal neurological deficits. CSF shows lymphocytic pleocytosis. Diagnosis: CSF PCR, MRI, EEG. Treatment: acyclovir for HSV, supportive care. Complications: permanent neurological deficits, cognitive impairment.

\medterm{Encephalocele}-- A neural tube defect where brain tissue and meninges herniate through a skull defect, most commonly in the occipital region. The herniated sac contains cerebrospinal fluid, meninges, and often brain tissue. Associated with other CNS anomalies including hydrocephalus and Chiari malformation. Presents as a visible, pulsatile midline mass. Treatment: surgical repair with reduction of herniated contents and dural closure. Prognosis depends on the amount of herniated brain tissue.

\medterm{Encephalopathy}-- Any diffuse disease of the brain that alters its structure or function, presenting with a global change in mental status (confusion, lethargy, personality change, coma) rather than focal neurologic deficits. Causes are numerous and often reversible: metabolic (hepatic encephalopathy--ammonia accumulation, uremic encephalopathy--uremic toxins, hypoglycemia, hyper/hyponatremia, hypercapnia, Wernicke encephalopathy--thiamine deficiency), toxic (drug overdose, alcohol intoxication/withdrawal, heavy metals), infectious (sepsis-associated encephalopathy, meningitis, encephalitis), hypoxic-ischemic (cardiac arrest), hypertensive (posterior reversible encephalopathy syndrome--PRES), and neurodegenerative (chronic traumatic encephalopathy, prion disease). Diagnosis: clinical assessment, metabolic panel, toxicology screen, neuroimaging (CT/MRI), EEG. Treatment: address underlying cause.

\medterm{Epidural Hematoma}-- Accumulation of blood between the skull and the dura mater, typically from arterial bleeding (middle meningeal artery) following temporal bone fracture. Presents with loss of consciousness, then a lucid interval (hours), followed by rapid neurological deterioration due to mass effect. Head CT: biconvex (lentiform), hyperdense collection that does not cross suture lines. Midline shift is common. Treatment: emergency burr hole or craniotomy for evacuation within 2-4 hours. Delay increases mortality. Outcome is excellent with prompt surgical intervention.

\medterm{Epilepsy}-- A neurological disorder characterized by recurrent, unprovoked seizures due to abnormal electrical activity in the brain. Classification includes focal, generalized, and unknown onset seizures. Etiology: genetic, structural, metabolic, infectious, immune. Diagnosis: clinical history, EEG, neuroimaging. Treatment: antiseizure medications, epilepsy surgery, vagus nerve stimulation, ketogenic diet for refractory cases.

\medterm{Epilepsy Surgery}-- Indicated for patients with drug-resistant epilepsy who have a concordant, resectable seizure focus identified on comprehensive presurgical evaluation. Failure of two or more appropriately chosen and adequately dosed antiepileptic drugs to achieve sustained seizure freedom. The presurgical evaluation includes scalp EEG with video monitoring, high-resolution MRI, PET scan (typically showing hypometabolism in the seizure focus), SPECT scan (ictal SPECT showing hyperperfusion and interictal SPECT showing hypoperfusion in the epileptic focus), neuropsychological assessment, and functional MRI for language and motor mapping. Invasive EEG monitoring using subdural grid electrodes or stereoelectroencephalography (SEEG) may be required when non-invasive data are discordant or the seizure focus is in or near eloquent cortex. The type of surgery depends on the location and nature of the epileptic focus: temporal lobectomy (the most common and most successful procedure, with 60-80 percent seizure freedom rates), extratemporal resection, hemispherotomy or hemispherectomy (for large hemispheric pathologies such as Rasmussen encephalitis or Sturge-Weber syndrome), corpus callosotomy (for drop attacks in generalised epilepsy), and vagus nerve stimulation (for drug-resistant epilepsy not amenable to resective surgery). Responsive neurostimulation (RNS) and deep brain stimulation of the anterior nucleus of the thalamus are newer neuromodulation options.

\medterm{Essential Tremor}-- The most common movement disorder, characterized by a bilateral, symmetrical, postural and kinetic tremor predominantly affecting the upper extremities, with a frequency of four to twelve hertz. It typically begins insidiously in the second or third decade of life and slowly progresses over years. The tremor is most apparent when the arms are held in a sustained posture (such as arms outstretched) and during intentional movement (such as reaching for an object), and it usually improves temporarily with alcohol consumption (a characteristic feature). The tremor can also affect the head (titubation), voice (vocal tremor, producing a quavering quality), legs, and trunk. Advanced cases may cause significant disability in writing, drinking from a cup, eating, and fine motor tasks. Essential tremor is often misdiagnosed as Parkinson disease, but key distinguishing features include the bilateral symmetrical presentation, predominant postural and kinetic tremor (rather than the resting tremor of Parkinson disease), lack of other parkinsonian signs such as bradykinesia and rigidity, and positive family history (autosomal dominant inheritance in up to 50 percent of cases, with incomplete penetrance). The pathophysiology is not fully understood but likely involves abnormal oscillatory activity in the cerebello-thalamo-cortical circuit, with the cerebellum playing a central role. Diagnosis is clinical, based on the Movement Disorder Society consensus criteria. Treatment includes beta-blockers (propranolol is first-line) and primidone as first-line medications, with topiramate, gabapentin, and benzodiazepines as second-line options. For medically refractory cases, botulinum toxin injections can reduce tremor amplitude, and deep brain stimulation of the ventral intermediate nucleus of the thalamus is highly effective. Focused ultrasound thalamotomy is a newer, incisionless ablative option.

\medterm{Facial paralysis}-- Loss of voluntary motor function of the facial muscles, most commonly from peripheral facial nerve (CN VII) dysfunction. The most frequent cause is Bell palsy (idiopathic peripheral facial paralysis), accounting for 60-70 percent of cases, with an incidence of 20-30 per 100,000 annually. Bell palsy presents as acute onset (over hours to 2-3 days) of unilateral lower motor neuron facial weakness involving both the upper and lower face (forehead, eye closure, mouth), often preceded by pain behind the ear (mastoid pain). Other causes of facial paralysis include Ramsay Hunt syndrome (herpes zoster oticus presenting with vesicles in the external auditory canal and severe pain), Lyme disease (erythema migrans and systemic symptoms), Guillain-Barre syndrome (bilateral facial weakness), otitis media, parotid tumors (slowly progressive), and stroke (upper motor neuron pattern sparing the forehead). Diagnosis involves careful history and physical examination to localise the lesion and identify the cause, with additional testing (MRI, nerve conduction studies, serology) as indicated. Management is directed at the underlying cause: Bell palsy is treated with corticosteroids, Ramsay Hunt syndrome with antivirals plus corticosteroids, Lyme disease with antibiotics, and tumors with surgical resection. Eye protection is critical in all cases of incomplete eye closure to prevent corneal abrasion and ulceration. Prognosis depends on the cause: Bell palsy has 70 percent complete recovery, while facial nerve tumors and trauma have variable outcomes.

\medterm{Friedreich Ataxia}-- The most common inherited ataxia, an autosomal recessive neurodegenerative disorder caused by a GAA trinucleotide repeat expansion in intron one of the frataxin gene (FXN) on chromosome 9q21.1. Normal individuals have five to thirty-three GAA repeats, while patients with Friedreich ataxia have typically sixty-six to one thousand seven hundred repeats on both alleles. The expanded repeat causes transcriptional silencing of the frataxin gene, leading to reduced frataxin protein, a mitochondrial protein involved in iron-sulphur cluster biogenesis and iron homeostasis. Reduced frataxin leads to mitochondrial iron accumulation, oxidative stress, and impaired mitochondrial function, particularly in neurons and cardiac myocytes. Clinical onset is typically between ages 10 and 15, with progressive gait and limb ataxia, dysarthria, loss of deep tendon reflexes (areflexia), loss of vibration and proprioception, extensor plantar responses (upper motor neuron signs due to corticospinal tract involvement), and scoliosis. Non-neurological features include hypertrophic cardiomyopathy (the most common cause of death), diabetes mellitus, and pes cavus. Diagnosis is clinical, supported by nerve conduction studies showing absent sensory nerve action potentials with preserved motor conduction velocities, and confirmed by genetic testing demonstrating biallelic GAA repeat expansions in the FXN gene. Management is supportive and multidisciplinary, including physical and occupational therapy, speech therapy, management of cardiomyopathy, treatment of diabetes, and orthopaedic management of scoliosis. The antioxidant idebenone has not shown consistent benefit in clinical trials. Prognosis is variable, with most patients becoming wheelchair-bound within 10-15 years of onset; mean survival is approximately 40-50 years, with cardiac failure as the leading cause of death.

\medterm{GABA (Gamma-Aminobutyric Acid)}-- The primary inhibitory neurotransmitter in the mammalian central nervous system, synthesised from glutamate via glutamic acid decarboxylase (GAD) with pyridoxal phosphate (vitamin B6) as cofactor. GABA acts on two classes of receptors: GABA-A (ionotropic, chloride ion channel, target of benzodiazepines, barbiturates, and ethanol) and GABA-B (metabotropic, G-protein coupled, modulates calcium and potassium channels). Functions: reduces neuronal excitability, regulates muscle tone, influences mood and anxiety. Implicated in: anxiety disorders (reduced GABAergic tone), epilepsy (GABA deficiency may contribute to seizures), Huntington disease (reduced GABA levels in basal ganglia), and spasticity (target of baclofen, a GABA-B agonist).

\medterm{Gait}-- The pattern or style of walking. Gait is a complex motor function involving coordination of the central and peripheral nervous systems, muscles, and joints. Gait assessment evaluates stance, stride length, cadence, arm swing, truncal stability, and turning. Abnormal gait patterns include: hemiplegic gait (circumduction, stroke), Parkinsonian gait (festinating, shuffling, reduced arm swing), ataxic gait (wide-based, unsteady, cerebellar or sensory; Romberg test positive in sensory ataxia), steppage gait (foot drop, e.g., common peroneal nerve palsy or Charcot-Marie-Tooth disease), waddling gait (hip girdle weakness, myopathy), antalgic gait (pain avoidance, shortened stance phase on affected side), and Trendelenburg gait (pelvic drop on the swing side, gluteus medius weakness).

\medterm{Gait Assessment}-- A crucial component of the neurological examination that provides significant diagnostic information about the underlying neurological condition. Normal gait requires intact motor, sensory, cerebellar, and vestibular systems, as well as intact cognition. Several characteristic gait patterns are associated with specific neurological disorders. The hemiplegic gait is seen in stroke and other upper motor neuron lesions, with the affected leg stiffly circumducted and the arm flexed and adducted at the side. The parkinsonian gait (festinating gait) is characterized by small, shuffling steps with reduced arm swing, flexed posture, and difficulty initiating or stopping movement, sometimes with freezing of gait. The sensory ataxic gait is broad-based with high-stepping (foot slap) and a positive Romberg sign, caused by impaired proprioception from peripheral neuropathy, tabes dorsalis, or posterior column disease. The cerebellar ataxic gait is broad-based, staggering, and lurching, with difficulty performing tandem gait, and is caused by cerebellar vermis pathology. The myopathic gait (waddling gait) features a broad-based stance, exaggerated lateral trunk sway, and Trendelenburg lurch due to weak hip girdle muscles, seen in muscular dystrophies. The apraxic gait (magnetic gait) is characterized by difficulty lifting the feet from the floor, as if glued to the ground, with initiation difficulty, seen in normal pressure hydrocephalus and frontal lobe lesions. Diagnosis of the gait pattern and the underlying neurological cause directs further investigation and management.

\medterm{Gargalesthesia}-- The sensation of tickling. A complex sensory experience mediated by light tactile stimulation of sensitive skin areas (axillae, soles, flanks, ribcage). Involves both spinal reflexes (goosebumps, withdrawal) and emotional processing in the anterior cingulate cortex and somatosensory cortex. Ticklishness requires an intact somatosensory system and a social/emotional context; self-tickling does not produce the sensation due to cerebellar cancellation of self-generated movement. Loss of gargalesthesia may occur in peripheral neuropathy or spinal cord lesions.

\medterm{Generalized Tonic-Clonic Seizure}-- A generalized tonic-clonic seizure (formerly grand mal seizure) is a seizure involving bilateral tonic stiffening followed by rhythmic clonic jerking of the extremities, with loss of consciousness. The tonic phase consists of sustained muscle contraction causing initial extension or flexion of the limbs, opisthotonus, and jaw clenching, lasting ten to twenty seconds. The clonic phase follows with rhythmic jerking of the limbs at a frequency of approximately two to four hertz, gradually slowing over time. The postictal phase follows with confusion, drowsiness, headache, and muscle soreness, lasting minutes to hours. During the generalised seizure, there is loss of consciousness, autonomic changes (tachycardia, hypertension, diaphoresis, pupillary dilation), and sometimes tongue biting, urinary incontinence, and cyanosis during the tonic phase. Generalised tonic-clonic seizures can be primary generalised (arising from both hemispheres simultaneously, often with a genetic basis) or secondary generalised (focal seizure spreading to become bilateral). Diagnosis is clinical, supported by EEG which may show generalised spike-wave discharges interictally. First-line treatment includes sodium valproate (for primary generalised) or levetiracetam, lamotrigine, and topiramate. Status epilepticus is a medical emergency defined as a generalised tonic-clonic seizure lasting longer than 5 minutes or recurrent seizures without return to baseline, requiring urgent management with benzodiazepines, followed by second-line anticonvulsants and anesthetic agents.

\medterm{Grand Mal}-- A type of generalised seizure characterized by bilateral synchronous epileptic activity involving the entire brain. Formerly called grand mal, now classified as bilateral tonic-clonic seizure according to the ILAE 2017 classification. Phases: (1) tonic phase (10-30 s): sudden loss of consciousness, sustained muscle contraction causing arching of the back (opisthotonos), extension of limbs, respiratory arrest with cyanosis; (2) clonic phase (30-60 s): rhythmic alternating contraction and relaxation of all muscle groups, initially rapid then slowing; (3) postictal phase (variable): confusion, headache, drowsiness, muscle soreness, sometimes Todd paralysis (transient focal neurological deficit). Complications: tongue biting, urinary incontinence, aspiration, injury from falls. Status epilepticus: continuous seizure activity >5 minutes or >2 seizures without full recovery. Emergency management: airway, benzodiazepines (lorazepam IV or midazolam IM), then levetiracetam/phenytoin/valproate if needed.

\medterm{Guillain-Barre Syndrome}-- An acute, immune-mediated polyneuropathy typically presenting with ascending paralysis following an antecedent infection (Campylobacter jejuni, CMV, EBV, Zika). Variants: acute inflammatory demyelinating polyneuropathy (AIDP, most common in North America), acute motor axonal neuropathy (AMAN, more common in Asia). Presents with symmetric ascending weakness, areflexia, and paresthesias. Respiratory failure occurs in 20-30\%. Diagnosis: CSF shows albuminocytologic dissociation (elevated protein, normal cell count); EMG/NCS. Treatment: IVIG or plasmapheresis (equally effective). Prognosis: most recover, but may have residual weakness.

\medterm{Gyrus}-- A ridge or convolution on the surface of the cerebral cortex, separated by sulci (grooves). The folded structure triples the surface area of the brain (approximately 2500 cm²) while fitting within the cranial vault. Major gyri: precentral gyrus (primary motor cortex, Brodmann area 4), postcentral gyrus (primary somatosensory cortex, areas 3,1,2), superior temporal gyrus (auditory cortex, area 41/42, includes Heschl gyrus), middle and inferior frontal gyri (prefrontal cortex, Broca area in dominant hemisphere), fusiform gyrus (face and object recognition), angular gyrus (reading and arithmetic), supramarginal gyrus (sensory integration), cingulate gyrus (limbic system, pain and emotion), hippocampal gyrus (parahippocampal, memory encoding). Gyral abnormalities may be seen in developmental disorders (polymicrogyria, pachygyria) and neurodegenerative diseases (gyral atrophy in Alzheimer disease).

\medterm{Headache}-- One of the most common neurologic symptoms, classified as primary (migraine, tension-type, cluster, and trigeminal autonomic cephalalgias) or secondary (attributed to an underlying structural, infectious, vascular, or metabolic cause). Tension-type headache is the most prevalent primary headache disorder, presenting as bilateral, pressing or tightening quality, mild to moderate intensity, lasting 30 minutes to 7 days, without nausea or photophobia/phonophobia. Migraine is the second most common primary headache disorder after tension-type headache. Secondary headaches have numerous causes, with red flags (SNOOP4: systemic symptoms, neurological signs, onset sudden, older age at onset, pattern change, prior headache history, precipitating factors, pregnancy, and postural aggravation) prompting investigation. Diagnosis is by careful history and neurological examination, with neuroimaging (CT or MRI) and lumbar puncture as indicated to exclude secondary causes. Management depends on the specific headache type: acute treatment includes simple analgesics, triptans, and antiemetics; preventive treatment includes beta-blockers, tricyclic antidepressants, antiepileptics, CGRP monoclonal antibodies, and nerve blocks.

\medterm{Headache and Migraine}-- Headaches are pain in the head or upper neck region, classified as primary (tension-type, migraine, cluster) or secondary to underlying conditions. Migraines are a neurological disorder characterized by recurrent episodes of moderate to severe throbbing pain, often accompanied by nausea, photophobia, and aura. Management includes acute treatments such as triptans and NSAIDs, plus preventive therapies including medications and lifestyle modifications.

\medterm{Hemifacial Spasm}-- Involuntary, unilateral, intermittent twitching of the facial muscles innervated by the facial nerve (CN VII). Typically begins around the eye (orbicularis oculi) and spreads to the lower face. Most commonly caused by vascular compression of the facial nerve at the root exit zone (anterior inferior cerebellar artery or posterior inferior cerebellar artery). Other causes: facial nerve injury, tumors, MS. Diagnosis: clinical; MRI/MRA to rule out structural cause. Treatment: botulinum toxin injections (first-line), microvascular decompression for refractory cases.

\medterm{Hemiplegia}-- Paralysis of one side of the body (the face, arm, and leg on the same side), caused by damage to the corticospinal tract (pyramidal tract) at any level from the cerebral cortex to the spinal cord. The most common cause of acute hemiplegia in adults is ischemic stroke affecting the contralateral cerebral hemisphere, particularly lesions of the internal capsule or motor cortex. Other causes include hemorrhagic stroke, brain tumors, traumatic brain injury, cerebral palsy (in children), and brain tumors. The clinical presentation includes contralateral facial weakness (lower face only, with forehead sparing due to bilateral cortical innervation of the upper face), contralateral arm and leg weakness (typically more severe in the arm and face than the leg in cortical strokes, or equal involvement in internal capsule or brainstem lesions), and associated neurological deficits depending on the location of the lesion (aphasia, neglect, visual field defects, sensory loss, dysarthria). The weakness is initially flaccid in the acute phase, with spasticity and hyperreflexia developing over days to weeks. Diagnosis is by neuroimaging (CT for acute hemorrhage, MRI for ischemic stroke) to identify the cause. Management is directed at the underlying cause: acute stroke management, rehabilitation, and prevention of secondary complications including contractures, pressure ulcers, DVT, and shoulder pain. Functional recovery depends on the extent of damage and the patient's age and pre-morbid status.

\medterm{Hemorrhagic Stroke}-- A neurological deficit caused by rupture of a cerebral blood vessel, accounting for 13\% of all strokes but 40\% of stroke mortality. Intracerebral hemorrhage: hypertensive (basal ganglia, thalamus, pons, cerebellum), amyloid angiopathy (lobar, elderly), vascular malformations, anticoagulation. Subarachnoid hemorrhage: aneurysm rupture. Presents with sudden severe headache, neurological deficits, decreased consciousness, and hypertension. Diagnosis: non-contrast CT is first-line. Management: ABCs, reversal of anticoagulation (vitamin K, FFP, PCC, idarucizumab for dabigatran, andexanet alfa for factor Xa inhibitors), blood pressure control (SBP <140-160), neurosurgical evacuation for cerebellar >3cm or lobar >30cc, external ventricular drain for hydrocephalus.

\medterm{Herpes Simplex Encephalitis}-- The most common cause of sporadic fatal encephalitis, caused by herpes simplex virus type 1 (HSV-1) in over ninety percent of cases. The virus typically infects the temporal lobes, orbitofrontal cortex, and insular cortex, causing hemorrhagic necrosis. Clinical presentation includes fever, headache, altered mental status, behavioral changes, seizures, focal neurological deficits (aphasia, hemiparesis), and personality changes. The temporal lobe predilection explains the common presenting features of personality changes, aphasia, and temporal lobe seizures. MRI brain is the imaging modality of choice, showing characteristic T2 hyperintensity and swelling of the temporal lobes (especially the medial temporal lobe including the hippocampus and amygdala), insular cortex, and orbitofrontal region, sometimes with areas of hemorrhage. The diagnosis is confirmed by detection of HSV-1 DNA in CSF by polymerase chain reaction (PCR), which has high sensitivity and specificity. CSF analysis typically shows a lymphocytic pleocytosis, elevated protein, and normal or slightly low glucose. EEG may show periodic lateralising epileptiform discharges (PLEDs) originating from the temporal region. Treatment with high-dose intravenous acyclovir (10 mg/kg every 8 hours for 14-21 days) should be started immediately when the diagnosis is suspected, even before confirmation by PCR, as early treatment significantly reduces mortality and improves neurological outcomes. Mortality is approximately 10-20 percent with treatment, and survivors often have long-term neurological deficits including memory impairment, personality changes, and epilepsy.

\medterm{Holoprosencephaly}-- A structural brain malformation resulting from incomplete or failed division of the forebrain (prosencephalon) into the cerebral hemispheres during the fifth week of gestation. Classified by severity: alobar (single ventricle, fused thalami, no interhemispheric fissure, most severe), semilobar (partial division), and lobar (nearly complete division, mildest). Presents with developmental delay, seizures, spasticity, and hypothalamic/pituitary dysfunction. Craniofacial anomalies correlate with severity: cyclopia, ethmocephaly, cebocephaly, or cleft lip/palate. Causes include trisomy 13 and mutations in SHH, ZIC2, SIX3, or TGIF genes. Treatment is supportive.

\medterm{Horner Syndrome}-- A clinical triad caused by disruption of the sympathetic pathway to the eye and face, resulting in ipsilateral miosis (constricted pupil), ptosis (drooping of the upper eyelid due to paralysis of the superior tarsal muscle), and anhidrosis (loss of sweating on the affected side of the face). The sympathetic pathway to the eye is a three-neuron chain: the first-order neuron descends from the hypothalamus through the brainstem to the ciliospinal center of Budge at C8-T2; the second-order neuron exits the spinal cord at C8-T2, ascends in the sympathetic chain, and synapses in the superior cervical ganglion; the third-order neuron travels along the internal carotid artery to innervate the pupillary dilator muscle, Muller muscle of the upper eyelid, and sweat glands of the face. Causes of Horner syndrome are classified by the site of the lesion along this three-neuron chain. Central (first-order) causes include brainstem stroke (Wallenberg syndrome), multiple sclerosis, syringobulbia, and spinal cord lesions above T1. Preganglionic (second-order) causes include Pancoast tumor (apical lung cancer), cervical rib, thoracic aortic aneurysm, thyroid carcinoma, neck trauma or surgery, and mediastinal masses. Postganglionic (third-order) causes include carotid artery dissection (a critical cause to identify), carotid aneurysm or thrombosis, cavernous sinus lesions, cluster headache, and otitis media. Diagnosis involves pharmacological testing with topical apraclonidine or cocaine drops, which confirm Horner syndrome, and topical hydroxyamphetamine to differentiate preganglionic from postganglionic lesions. Urgent imaging of the brain, neck, and chest (CT, MRI, CTA, or MRA) is required to identify the underlying cause, with carotid dissection being a medical emergency requiring antithrombotic therapy.

\medterm{Huntington Disease}-- An autosomal dominant neurodegenerative disorder caused by a CAG trinucleotide repeat expansion in the huntingtin (HTT) gene on chromosome 4p16.3. The normal number of CAG repeats is ten to thirty-five, while forty or more repeats are fully penetrant, causing the disease. There is an inverse correlation between the number of CAG repeats and the age of onset: larger repeat expansions cause earlier onset. The disease is characterized by a classic triad of motor, cognitive, and psychiatric abnormalities. The motor features include chorea (rapid, involuntary, dance-like movements) as the most characteristic feature, which initially may be subtle—mild fidgetiness, clumsiness, or involuntary facial movements—and progresses to involve the trunk and limbs, affecting gait, speech, and swallowing. Dystonia and parkinsonism can also occur, particularly in juvenile-onset cases. Cognitive decline begins with executive dysfunction (reduced flexibility, poor planning, and impulsivity), progressing to global dementia. Psychiatric manifestations are common and include depression (the most common), apathy, irritability, obsessive-compulsive symptoms, and psychosis. Diagnosis is clinical, confirmed by genetic testing demonstrating 40 or more CAG repeats in the HTT gene. Presymptomatic testing is available for individuals at risk, requiring careful genetic counseling. MRI brain shows atrophy of the caudate nucleus and putamen (striatal atrophy) with corresponding enlargement of the frontal horns of the lateral ventricles. There is no disease-modifying treatment; management is symptomatic with antipsychotics (particularly tetrabenazine and deutetrabenazine for chorea), antidepressants, and anxiolytics. Mean survival is 15-20 years from symptom onset.

\medterm{Huntington's Disease (Chorea)}-- An autosomal dominant neurodegenerative disorder caused by a CAG trinucleotide repeat expansion (>36 repeats, normally <26) in the HTT gene on chromosome 4 (4p16.3), which encodes the huntingtin protein. The expanded repeat produces an elongated polyglutamine tract causing toxic protein aggregation and neuronal death, especially in the caudate nucleus, putamen, and cerebral cortex. Anticipation: repeat length tends to increase in successive generations, especially with paternal transmission (daughters of affected fathers are more likely to develop earlier onset). Prevalence: 5-10 per 100,000. Age of onset: typically 35-50 years; juvenile onset (<20 years, large repeats >60 CAG) presents with rigidity and seizures. Clinical triad: (1) Movement disorder: chorea (irregular, rapid, jerky, purposeless movements of face, trunk, limbs; initially subtle may be masked by voluntary movement), dystonia, parkinsonism (rigidity, bradykinesia, especially in later stages), impaired voluntary movement (incoordination, dysarthria, dysphagia). (2) Cognitive decline: executive dysfunction (impaired planning, organising, multitasking, cognitive flexibility), memory deficits, psychomotor slowing, progressing to subcortical dementia. (3) Psychiatric symptoms: depression (most common, 30-50\%, may precede motor onset), irritability, apathy, anxiety, obsessive-compulsive symptoms, psychosis (hallucinations, delusions, 5-15\%). Diagnosis: clinical (positive family history in 90\%, diagnostic criteria: motor symptoms in at-risk individual), genetic testing (CAG repeat length, confirmatory). Presymptomatic testing offered with genetic counseling. Imaging: MRI/CT shows caudate atrophy (bicaudate ratio increased) and ex vacuo ventricular enlargement. Management: no disease-modifying treatment. Chorea: tetrabenazine (VMAT2 inhibitor, depletes dopamine), deutetrabenazine, valbenazine; antipsychotics (olanzapine, risperidone) for chorea and psychiatric symptoms. Depression: SSRIs. Multidisciplinary care: physiotherapy, speech therapy, occupational therapy, nutritional support (high-calorie diet). Disease progression: inexorable, median survival 15-20 years from onset, death typically from pneumonia or complications of immobility. Research: antisense oligonucleotides (ASO) targeting HTT mRNA are in clinical trials.

\medterm{Hydranencephaly}-- A rare congenital condition where the cerebral hemispheres are nearly completely absent, replaced by thin-walled, fluid-filled sacs, while the brainstem, basal ganglia, and cerebellum are relatively preserved. Caused by bilateral internal carotid artery occlusion in utero (vascular insult, infection, or twin-twin transfusion). Presents with normal head size at birth, rapid head growth (mimicking hydrocephalus), preserved brainstem reflexes (sucking, swallowing, eye movements), but no cognitive development, visual tracking, or voluntary movement. Diagnosis by transillumination (glows brightly) and MRI/CT showing absent cerebral tissue with preserved posterior fossa structures. Differentiated from severe hydrocephalus (which has a thin cortical rim). No treatment; supportive care. Most die in infancy or childhood.

\medterm{Hydrocephalus}-- Abnormal accumulation of cerebrospinal fluid (CSF) within the ventricles of the brain, leading to increased intracranial pressure. Communicating hydrocephalus: impaired CSF absorption (subarachnoid hemorrhage, meningitis, leptomeningeal carcinomatosis). Non-communicating (obstructive) hydrocephalus: blockage within the ventricular system (aqueductal stenosis, mass lesions, colloid cyst). Normal pressure hydrocephalus (NPH): classic triad of gait apraxia, urinary incontinence, and cognitive decline (dementia); enlarged ventricles on imaging with normal opening pressure. In infants: rapidly enlarging head circumference, bulging fontanelles, sunsetting eyes. Treatment: ventriculoperitoneal shunt (mainstay), endoscopic third ventriculostomy for selected obstructive cases.

\medterm{Insomnia}-- A sleep disorder characterized by difficulty falling asleep, staying asleep, or waking too early despite adequate opportunity for sleep, with resulting daytime impairment. It may be acute or chronic and is often triggered by stress, anxiety, depression, or poor sleep habits. Treatment includes cognitive behavioral therapy for insomnia (CBT-I) as first-line therapy, with medications considered for short-term use.

\medterm{Internuclear Ophthalmoplegia}-- A conjugate horizontal gaze disorder caused by a lesion of the medial longitudinal fasciculus (MLF), a heavily myelinated fiber tract that connects the abducens nucleus (CN VI) on one side to the oculomotor nucleus (CN III) on the contralateral side, coordinating horizontal gaze. In a left INO (the most common type), a lesion of the left MLF impairs adduction of the left eye (the eye cannot move past the midline) and causes abducting nystagmus of the right eye (the abducting eye develops jerky nystagmus as it tries to overcome the impaired adduction signal). Convergence is preserved because the medial rectus subnucleus in CN III is also innervated via the vergence pathway, which bypasses the MLF. INO is most commonly bilateral in multiple sclerosis and unilateral in ischemic stroke affecting the brainstem (particularly the pons). The lesion is in the ipsilateral MLF relative to the eye that fails to adduct. Diagnosis is clinical by examination of horizontal gaze in both directions, with the characteristic pattern of impaired adduction on the ipsilateral side and abducting nystagmus on the contralateral side. MRI brain demonstrates the MLF lesion and any associated demyelinating or ischemic pathology. Treatment is directed at the underlying cause, with symptomatic management of oscillopsia from the nystagmus (gabapentin, memantine, clonazepam).

\medterm{Intracerebral Hemorrhage}-- Bleeding directly into the brain parenchyma, accounting for approximately ten to fifteen percent of all strokes and carrying a high mortality rate of approximately forty percent at thirty days. The most common cause is chronic hypertension, which causes lipohyalinosis and Charcot-Buchard microaneurysm formation in the penetrating arterioles of the basal ganglia (most commonly the putamen), thalamus, pons, and cerebellum. Other causes include cerebral amyloid angiopathy (common cause of lobar ICH in elderly patients), vascular malformations (arteriovenous malformations, cavernous malformations), intracranial aneurysms, coagulopathy (warfarin, DOACs, thrombocytopenia), brain tumors (tumoural hemorrhage), vasculitis, cerebral venous sinus thrombosis with venous infarction, and drug abuse (cocaine, amphetamines). Clinical presentation includes sudden onset of severe headache, vomiting, focal neurological deficit (hemiparesis, aphasia, sensory loss, visual field defects), altered level of consciousness, and seizures. The deficit progresses over minutes to hours as the haematoma expands. Diagnosis is by non-contrast CT head, which shows a hyperdense (white) lesion within the brain parenchyma. MRI is more sensitive for detecting the underlying cause (amyloid angiopathy, cavernous malformations, tumors). CT angiography may detect a spot sign indicating active extravasation and risk of haematoma expansion. Management includes blood pressure control (systolic <140 mmHg to reduce haematoma expansion), reversal of anticoagulation (vitamin K, prothrombin complex concentrate, protamine), intracranial pressure monitoring and management, and surgical haematoma evacuation for cerebellar hemorrhage greater than 3 cm, lobar hemorrhage with significant mass effect, or superficial hemorrhage in a deteriorating patient. The STICH trial showed no overall benefit for early surgery in spontaneous supratentorial ICH.

\medterm{Intracranial Hemorrhage}-- Bleeding within the cranial cavity. Epidural hematoma: arterial bleeding (middle meningeal artery), typically from temporal bone fracture; lucid interval then rapid deterioration. Subdural hematoma: venous bleeding from bridging veins; acute (<72h), subacute, or chronic (elderly, alcoholics). Subarachnoid hemorrhage: arterial bleeding (aneurysm rupture 85\%); thunderclap headache, nuchal rigidity. Intracerebral hemorrhage: hypertensive (basal ganglia, thalamus, pons, cerebellum), amyloid angiopathy, vascular malformations. Intraventricular hemorrhage: extension from ICH or primary. Diagnosis: non-contrast CT is first-line. Treatment: ABCs, reversal of anticoagulation, blood pressure control (SBP <140 for ICH), surgical evacuation (cerebellar >3cm, lobar >30cc), external ventricular drain for hydrocephalus, aneurysm clipping/coiling for SAH.

\medterm{Ischemic Stroke}-- A neurological deficit caused by focal cerebral ischemia due to thrombotic or embolic occlusion of a cerebral artery (87\% of all strokes). Thrombotic: in situ occlusion of atherosclerotic vessel. Embolic: embolus from cardiac source (atrial fibrillation, valvular disease, patent foramen ovale) or large artery (carotid/vertebral). Lacunar: small vessel occlusion due to lipohyalinosis. Presents with sudden focal deficits (hemiparesis, aphasia, neglect, ataxia) corresponding to vascular territory. Diagnosis: non-contrast CT (rule out hemorrhage), CT perfusion, MRI DWI (gold standard). Acute management: IV tPA within 4.5 hours, mechanical thrombectomy for large vessel occlusion within 24 hours. Secondary prevention: antiplatelets, statins, blood pressure control, anticoagulation for cardioembolic.

\medterm{Jacksonian Epilepsy (Focal Seizures)}-- A type of focal (partial) seizure originating in the primary motor cortex, characterized by clonic (rhythmic jerking) or tonic (sustained contraction) movements that spread sequentially to adjacent body parts along the motor homunculus (Jacksonian march). Named after John Hughlings Jackson (1835-1911), who first described the phenomenon. The march typically begins in the hand or face (largest cortical representation) and spreads proximally up the arm, then to the leg on the same side. Can also present as sensory seizures (somatosensory cortex, paraesthesiae) or versive seizures (frontal eye field). Consciousness is usually preserved during the seizure (focal aware seizure, ILAE 2017 classification). May evolve into a bilateral tonic-clonic seizure (secondary generalisation). Aetiologies: any focal brain lesion--cortical dysplasia, tumor (gliomas, meningiomas), stroke, arteriovenous malformation, traumatic brain injury, infection (abscess, encephalitis), and mesial temporal sclerosis. Diagnosis: EEG (may show focal spikes or sharp waves over the contralateral central region), MRI brain with epilepsy protocol. Treatment: antiseizure medications (carbamazepine, lamotrigine, levetiracetam, oxcarbazepine). Surgical resection of the epileptogenic focus may be curative in 60-80\% of drug-resistant cases.

\medterm{Karnofsky Performance Status}-- The Karnofsky Performance Status (KPS) scale is a widely used clinical tool for assessing the functional status of patients, particularly those with cancer, and quantifies the patient's ability to perform activities of daily living and tolerate treatments. The scale ranges from zero (dead) to one hundred (normal, no complaints, no evidence of disease) in increments of ten. KPS scores of eighty to one hundred indicate that the patient is able to carry on normal activity and work, with no special care needed; KPS 70 indicates that the patient is able to care for self but cannot carry on normal activity or do active work; KPS 50 indicates that the patient requires considerable assistance and frequent medical care; and KPS 20 indicates that the patient is very sick, requiring active supportive treatment. The scale was developed by Karnofsky and Burchenal in 1949 and is used in oncology to determine prognosis, guide treatment intensity, and assess eligibility for clinical trials. The KPS has been validated as an independent predictor of survival in many cancers. A score of 70 or higher is typically required for aggressive treatments such as combination chemotherapy or major surgery, while lower scores indicate the need for palliative and supportive care. The scale has limitations, including inter-rater variability and a lack of sensitivity to small changes in function.

\medterm{Limbic System}-- A set of interconnected brain structures located in the medial temporal lobe and around the corpus callosum, involved in emotion, memory, motivation, behaviour, and autonomic regulation. The term was coined by Paul Broca (1878, le grand lobe limbique) and expanded by Papez (1937, Papez circuit) and MacLean (1952, 'visceral brain'). Key structures: (1) Hippocampus: memory formation (spatial, episodic), encoding and consolidation from short-term to long-term memory. (2) Amygdala: emotional processing (fear, anger, anxiety, pleasure), emotional memory, social behaviour, and threat detection. (3) Cingulate gyrus: emotional processing, conflict monitoring, pain processing, decision-making. (4) Hypothalamus: autonomic (sympathetic and parasympathetic), neuroendocrine (CRH, TRH, GHRH, GnRH, ADH), hunger/satiety, thirst, thermoregulation, circadian rhythms. (5) Fornix: major output tract from hippocampus to mammillary bodies. (6) Mammillary bodies: memory (reciprocal connections with hippocampus via fornix, and with thalamus). (7) Anterior thalamic nucleus: part of the Papez circuit, spatial memory. (8) Septal nuclei: reward, pleasure, reinforcement. (9) Parahippocampal gyrus: input to hippocampus. (10) Olfactory tubercle: smell. Limbic system dysfunction causes: memory disorders (Alzheimer disease--hippocampal atrophy, Korsakoff syndrome--mammillary body and thalamic damage), emotional dysregulation (anxiety, depression, PTSD, phobias), seizure disorders (temporal lobe epilepsy), schizophrenia (disrupted limbic connectivity), and autonomic dysfunction.

\medterm{Lissencephaly}-- A severe brain malformation characterized by a smooth cerebral surface due to absent (agyria) or decreased (pachygyria) sulcation and gyration, caused by impaired neuronal migration during 12-16 weeks of gestation. Presents with severe developmental delay, intractable epilepsy (infantile spasms, Lennox-Gastaut), spasticity, and feeding difficulties. Facial features may include bitemporal hollowing, prominent forehead, and micrognathia. Diagnosis by MRI showing a smooth brain surface, thick cortex, and shallow or absent Sylvian fissures. Genetic causes include LIS1 (Miller-Dieker syndrome, 17p13.3 deletion) and DCX (X-linked) mutations. Prognosis is poor; most children have severe disabilities and shortened lifespan.

\medterm{Locked-in Syndrome}-- A neurological condition in which the patient is fully conscious and aware but cannot move or speak due to complete paralysis of all voluntary muscles except for vertical eye movements and blinking (total locked-in syndrome: loss of all eye movements). The patient is 'locked in' a body that cannot respond but has intact cognition, sensation, and consciousness. Most common cause: bilateral ventral pontine stroke (basilar artery occlusion/ thrombosis, infarction of the ventral pons, damaging the corticospinal tracts and corticobulbar tracts, sparing the reticular activating system (consciousness) and the midbrain tectum (vertical eye movements). Other causes: central pontine myelinolysis (rapid correction of hyponatraemia), pontine hemorrhage, brainstem tumor, encephalitis (enterovirus 71), trauma (basilar skull fracture), amyotrophic lateral sclerosis (advanced), and severe Guillain-Barré syndrome. Diagnosis: clinical (consciousness confirmed by eye movement communication). MRI brain shows ventral pontine lesion. EEG is normal (distinguishes from vegetative state/minimally conscious state). Communication: often established through eye movements (blinking code, eye-tracker technology, EEG-based brain-computer interfaces). Prognosis: poor for functional recovery; mortality 60-90\% in the first year. Some patients survive many years with intensive support and communication devices. The condition is distinguished from: coma (no eye opening, no awareness), vegetative state (wakefulness without awareness, no purposeful responses), and minimally conscious state (inconsistent but reproducible evidence of awareness).

\medterm{Locomotor Ataxia (Tabes Dorsalis)}-- A late manifestation of untreated tertiary syphilis (neurosyphilis) caused by infection with Treponema pallidum, involving progressive degeneration of the dorsal columns (fasciculus gracilis and cuneatus) and dorsal roots of the spinal cord. The demyelination and neuronal loss affect the sensory pathways mediating proprioception (joint position sense) and vibration sense. Clinical features: sensory ataxia (wide-based, stamping gait, worse with eyes closed, positive Romberg sign), paraesthesiae (lightning pains--sudden, severe, stabbing pain in the legs, back, trunk), loss of deep tendon reflexes (especially knee and ankle jerks, areflexia due to dorsal root involvement), impaired proprioception and vibration sense, Charcot joints (neuropathic arthropathy, painless joint destruction, most commonly knee, ankle), Argyll Robertson pupils (bilateral small irregular pupils that accommodate but do not react to light, light-near dissociation), and optic atrophy. Onset is typically 15-30 years after primary infection. Diagnosis: clinical, serum treponemal tests (TPPA, FTA-ABS) positive, CSF analysis (lymphocytic pleocytosis, elevated protein, positive CSF VDRL), and spinal cord MRI (dorsal column T2 hyperintensity and atrophy). Treatment: aqueous crystalline penicillin G 18-24 million units IV daily for 10-14 days. Prognosis: some improvement in symptoms, but neurological deficits are often permanent.

\medterm{Lumbar Disc Herniation}-- The most common cause of lumbar radiculopathy (sciatica), occurring when the nucleus pulposus of a lumbar intervertebral disc herniates through the annulus fibrosus, compressing adjacent nerve roots. The most commonly affected level is L4-L5 and L5-S1, accounting for approximately ninety-five percent of lumbar disc herniations. L4-L5 disc herniation typically compresses the L5 nerve root, causing pain radiating from the buttock down the lateral leg to the dorsum of the foot and great toe, weakness of ankle dorsiflexors and evertors (foot drop), and no reflex changes. L5-S1 disc herniation typically compresses the S1 nerve root, causing pain radiating from the buttock down the posterior leg to the heel and lateral foot, weakness of plantar flexion (gastrocnemius, soleus) and ankle eversion, and decreased ankle jerk reflex. Cauda equina syndrome is a surgical emergency requiring urgent MRI and neurosurgical referral. The majority of lumbar disc herniations resolve with conservative management including NSAIDs, activity modification, and physiotherapy. Indications for surgical intervention include cauda equina syndrome, progressive neurological deficit, and persistent radicular pain refractory to 6-12 weeks of conservative management. Surgical options include microdiscectomy, endoscopic discectomy, and minimally invasive tubular discectomy.

\medterm{Lumbar puncture}-- Also known as spinal tap, a procedure to obtain cerebrospinal fluid from the subarachnoid space in the lumbar spine, typically at the L3-L4 or L4-L5 interspace (below the termination of the spinal cord at L1-L2 in adults). Indications include diagnosis of meningitis and encephalitis (CSF Gram stain, culture, cell count, glucose, protein, PCR for viruses and bacteria), subarachnoid hemorrhage (xanthochromia from bilirubin breakdown products, RBC count not clearing from tube 1 to tube 4), neuroinflammatory and demyelinating diseases (multiple sclerosis, Guillain-Barre syndrome, chronic inflammatory demyelinating polyneuropathy), and neoplastic conditions (leptomeningeal carcinomatosis, lymphoma). The procedure involves inserting a hollow needle (typically 22G or 20G) into the subarachnoid space at the L3-L4 or L4-L5 interspace, after sterile preparation and local anesthesia. The patient is positioned in the lateral decubitus (fetal) position or sitting and leaning forward to widen the interlaminar space. Opening pressure is measured using a manometer before CSF is collected, with normal opening pressure ranging from 10-20 cm H2O. CSF is collected in sequential tubes for cell count and differential, glucose and protein, Gram stain and culture, and additional studies (PCR, cytology, oligoclonal bands, flow cytometry) as indicated. Contraindications include increased intracranial pressure with mass effect or intracranial mass (risk of brain herniation—the most feared complication), coagulopathy or thrombocytopenia, infection at the puncture site, and spinal cord compression. Complications include post-dural puncture headache (the most common complication, caused by CSF leak through the dural puncture, managed with bed rest, caffeine, hydration, and epidural blood patch for severe cases), bleeding (spinal epidural haematoma, rare but serious), infection (meningitis, epidural abscess), and radicular pain from nerve root irritation.

\medterm{Lumbar Radiculopathy}-- A condition caused by compression, irritation, or inflammation of a spinal nerve root in the lumbar spine, resulting in pain radiating into the lower extremity along the distribution of the affected nerve root (sciatica). The most commonly affected nerve roots are L4, L5, and S1. L4 radiculopathy causes pain radiating into the medial leg and medial foot, weakness of knee extension (quadriceps), and decreased knee reflex. L5 radiculopathy causes pain radiating into the lateral leg and dorsum of the foot, weakness of ankle dorsiflexion (foot drop) and great toe extension (extensor hallucis longus), with no reflex changes. S1 radiculopathy causes pain radiating into the posterior leg and lateral foot, weakness of plantar flexion (gastrocnemius, soleus), and decreased ankle jerk reflex. Diagnosis is clinical, confirmed by MRI of the lumbar spine demonstrating nerve root compression from disc herniation, spinal stenosis, or foraminal narrowing. EMG/NCS may confirm the affected nerve root and exclude peripheral neuropathy or plexopathy. Management is initially conservative with NSAIDs, activity modification, and physiotherapy. Epidural corticosteroid injections provide temporary relief. Surgical intervention (microdiscectomy or decompressive laminectomy) is indicated for cauda equina syndrome, progressive motor deficit, or persistent radicular pain despite 6-12 weeks of conservative management.

\medterm{Memory}-- The cognitive process of encoding, storing, and retrieving information, mediated by distributed neural networks involving the hippocampus, medial temporal lobe, prefrontal cortex, and limbic system. Types: sensory memory (brief, millisecond--second), short-term/working memory (seconds--minutes, limited capacity ~7 items), and long-term memory (subdivided into explicit/episodic and semantic and implicit/procedural). Memory impairment (amnesia) can be acute (transient global amnesia, concussion, seizure, medication side effects) or progressive (Alzheimer disease, vascular dementia, frontotemporal dementia, Lewy body dementia). Evaluation includes mental status examination (MMSE, MoCA), collateral history, laboratory workup (B12, TSH, syphilis, HIV), and brain imaging (MRI for atrophy pattern, vascular disease). Management: treat reversible causes, cognitive stimulation, safety planning, and pharmacotherapy (cholinesterase inhibitors for Alzheimer disease; memantine for moderate--severe stages).

\medterm{Meningioma}-- The most common primary intracranial tumor, accounting for approximately thirty-seven percent of all primary brain tumors, and arises from arachnoid cap cells of the meninges. Most meningiomas are benign (WHO grade I), slow-growing, and well-circumscribed, with a whorled histopathological pattern and psammoma bodies. Atypical (grade II) and anaplastic (grade III) meningiomas are more aggressive with higher recurrence rates. Meningiomas are more common in women, with a female-to-male ratio of approximately 2-3:1, and their incidence increases with age. Most meningiomas are asymptomatic and discovered incidentally on neuroimaging. Symptomatic meningiomas present with symptoms related to mass effect on adjacent structures, including headache, seizures (focal or generalised), and focal neurological deficits. The most common locations are parasagittal/falcine region, cerebral convexities, sphenoid wing, and olfactory groove. Sphenoid wing meningiomas may present with proptosis and visual disturbances. Diagnosis is by MRI with gadolinium, which shows a dural-based, extra-axial mass with homogeneous contrast enhancement and a dural tail sign. CT may show hyperostosis of the adjacent bone. Asymptomatic or small meningiomas with no mass effect are managed with observation and serial imaging. Symptomatic, growing, or large tumors are treated with surgical resection (Simpson grade I-V, with grade I meaning complete removal including dural attachment). Adjuvant radiotherapy or stereotactic radiosurgery is used for residual, recurrent, atypical, or malignant meningiomas. WHO grade I meningiomas have excellent prognosis, with 5-year recurrence rates of 5-10 percent after complete resection; grade II (atypical) have 30-40 percent recurrence; and grade III (anaplastic/malignant) have 50-80 percent recurrence.

\medterm{Meningitis}-- Inflammation of the meninges most commonly caused by bacterial infection (Neisseria meningitidis, Streptococcus pneumoniae, Haemophilus influenzae, Group B Streptococcus). Presents with fever, nuchal rigidity, photophobia, and altered mental status. Diagnosis: lumbar puncture with CSF analysis. Bacterial meningitis is a medical emergency requiring prompt empiric antibiotics and dexamethasone. Vaccination prevents several causes.

\medterm{Migraine}-- A primary headache disorder characterized by recurrent episodes of moderate to severe, typically unilateral, pulsating headache lasting 4-72 hours, associated with nausea, vomiting, photophobia, and phonophobia. Migraine is the second most disabling condition worldwide (second only to low back pain in years lived with disability), affecting 15-18 percent of women and 6-8 percent of men globally. The pathophysiology involves a complex neurovascular cascade: cortical spreading depression (a wave of neuronal depolarisation followed by prolonged suppression of cortical electrical activity that propagates across the cortex at 2-6 mm per minute), which is thought to be the pathophysiological substrate of the migraine aura, followed by activation of the trigeminovascular system (trigeminal nerve endings surrounding intracranial blood vessels release vasoactive neuropeptides such as CGRP, substance P, and neurokinin A), causing neurogenic inflammation and vasodilation of intracranial and extracranial blood vessels, which activates nociceptive pathways to the thalamus and cortex. Triggers include stress (the most common trigger), hormonal fluctuations (menstrual migraine), lack of sleep or excessive sleep, skipped meals, dehydration, alcohol (especially red wine), caffeine withdrawal, bright lights, loud noises, strong odours, and certain foods. Migraine is classified as episodic (fewer than 15 headache days per month) or chronic (15 or more headache days per month for more than 3 months, with at least 8 days meeting migraine criteria). Chronic migraine often evolves from episodic migraine and is associated with overuse of acute medications. Management of acute attacks includes simple analgesics (ibuprofen, paracetamol, aspirin), triptans (sumatriptan, rizatriptan, eletriptan—serotonin 5-HT1B/1D receptor agonists that inhibit CGRP release and cause vasoconstriction), antiemetics (metoclopramide, domperidone), and gepants (ubrogepant, rimegepant—CGRP receptor antagonists). Preventive therapy is indicated for patients with frequent or disabling attacks and includes beta-blockers (propranolol, metoprolol), tricyclic antidepressants (amitriptyline), antiepileptics (topiramate, sodium valproate), CGRP monoclonal antibodies (erenumab, galcanezumab, fremanezumab, eptinezumab), and botulinum toxin type A (OnabotulinumtoxinA) for chronic migraine.

\medterm{Migraine With Aura}-- A primary headache disorder characterized by recurrent attacks of fully reversible visual, sensory, or speech symptoms that develop gradually over at least five minutes and last no more than sixty minutes, followed by headache with the characteristics of migraine. Visual aura is the most common type, affecting approximately ninety percent of patients with migraine with aura, and typically consists of a fortification spectrum (scintillating scotoma) that expands over the visual field over 20-60 minutes, leaving a scotoma (blind spot) in its wake. Sensory aura is the second most common type, presenting with unilateral tingling and numbness that spreads slowly from the hand up the arm to the face, typically lasting 5-30 minutes. Speech and language aura is less common and presents with dysphasic speech. The migraine aura is thought to be caused by cortical spreading depression (CSD)—a wave of neuronal and glial depolarisation followed by prolonged suppression of neural activity that propagates across the cerebral cortex at a rate of 2-6 mm per minute. The visual aura localises to the occipital cortex, sensory aura to the somatosensory cortex, and speech aura to the language areas. If the aura occurs without headache (migraine aura without headache, formerly called acephalgic migraine), other causes such as transient ischemic attack and seizure must be excluded, particularly in older patients. Triptans and NSAIDs are effective for acute treatment, and they are safe to take during the aura phase. Preventive treatment is indicated for frequent or disabling attacks (typically more than 4 per month) and includes the same options as for migraine without aura. The risk of ischemic stroke is slightly elevated in migraine with aura, particularly in women who smoke and use oral contraceptives, and this should be considered in management.

\medterm{Migraine Without Aura}-- A primary headache disorder characterized by recurrent attacks of moderate to severe, pulsating, unilateral headache lasting four to seventy-two hours, associated with nausea, vomiting, photophobia, and phonophobia. The International Classification of Headache Disorders (ICHD-3) diagnostic criteria require at least five attacks fulfilling the following: headache lasting four to seventy-two hours, with at least two of the following four characteristics: unilateral location, pulsating quality, moderate or severe pain intensity, and aggravation by or causing avoidance of routine physical activity, and during the headache attack at least one of nausea and/or vomiting or photophobia and phonophobia. The attack is not better accounted for by another diagnosis. Migraine without aura is the most common subtype of migraine, accounting for approximately 70-80 percent of all migraine attacks. The headache phase is thought to be mediated by activation of the trigeminovascular system, with release of vasoactive neuropeptides such as CGRP causing neurogenic inflammation of the meningeal blood vessels. The pathophysiology of migraine without aura is less well understood than migraine with aura, as cortical spreading depression is not thought to occur. Management of acute attacks includes stratified care based on headache severity: mild attacks are treated with simple analgesics (ibuprofen, paracetamol, aspirin, diclofenac) plus an antiemetic; moderate-to-severe attacks are treated with triptans (sumatriptan, rizatriptan, eletriptan) or gepants (rimegepant, ubrogepant). Triptans are most effective when taken early during the headache phase before central sensitisation develops. Preventive therapy is indicated for patients with frequent or disabling attacks and includes beta-blockers (propranolol, metoprolol), tricyclic antidepressants (amitriptyline), antiepileptics (topiramate, sodium valproate), and CGRP monoclonal antibodies.

\medterm{Mild Cognitive Impairment}-- A transitional cognitive state between normal aging and dementia, characterized by subjective cognitive complaints confirmed by objective cognitive impairment that is not severe enough to significantly interfere with daily activities. MCI is classified as amnestic MCI (aMCI), where memory is the primary domain affected, and non-amnestic MCI (naMCI), where other cognitive domains (attention, executive function, visuospatial ability, language) are primarily affected. The amnestic subtype is the most common and is considered a prodromal stage of Alzheimer disease in most cases. The prevalence of MCI in adults over 60 is approximately 15-20 percent, and the annual conversion rate from MCI to dementia is approximately 10-15 percent, although some patients remain stable or revert to normal cognition. Diagnosis requires objective cognitive testing confirming impairment one to two standard deviations below age- and education-matched norms, with intact functional abilities. Biomarkers such as CSF A-beta and tau levels, amyloid PET imaging, and FDG-PET can help determine the likelihood of progression to Alzheimer disease. Management includes addressing modifiable vascular risk factors (hypertension, diabetes, smoking), cognitive stimulation, physical exercise, and monitoring for cognitive decline. There is no pharmacological treatment approved for MCI; acetylcholinesterase inhibitors are not recommended as they do not prevent progression to dementia. Prognosis is variable, with some patients remaining stable indefinitely while others progress to dementia within a few years.

\medterm{Multiple Sclerosis}-- A chronic autoimmune demyelinating disease of the central nervous system, most commonly presenting in young adults. Characterized by relapsing-remitting episodes of neurological dysfunction (optic neuritis, transverse myelitis, brainstem syndromes). Diagnosis: McDonald criteria with MRI dissemination in space and time. Disease-modifying therapies reduce relapse rate and disability progression.

\medterm{Multiple System Atrophy}-- A progressive, fatal neurodegenerative disorder characterized by varying combinations of autonomic failure, cerebellar ataxia, parkinsonism, and pyramidal signs. It is classified into two subtypes: MSA-P (Parkinsonian variant, previously striatonigral degeneration), characterized by prominent parkinsonism with poor levodopa response, and MSA-C (cerebellar variant, previously sporadic olivopontocerebellar atrophy), characterized by prominent cerebellar ataxia. Pathologically, MSA is characterized by the accumulation of alpha-synuclein in oligodendrocytes, forming glial cytoplasmic inclusions (GCIs), which is the pathological hallmark and distinguishes MSA from other synucleinopathies such as Parkinson disease and Lewy body dementia where alpha-synuclein accumulates in neurons. Onset is typically in the 50s or 60s, with rapidly progressive course leading to severe disability within 5-10 years. Clinical features of autonomic failure (present in both subtypes) include orthostatic hypotension (symptomatic drop in blood pressure on standing, often the earliest symptom), urinary incontinence, erectile dysfunction, and impaired sweating. Other features include stridor (nocturnal and diurnal, indicating vocal cord abductor paralysis, a potentially life-threatening complication requiring CPAP or tracheostomy), REM sleep behaviour disorder, and respiratory dysfunction. The parkinsonism in MSA-P is typically symmetrical, with poor or absent response to levodopa, early postural instability, and craniocervical dystonia. The cerebellar features in MSA-C include gait and limb ataxia, dysarthria, and oculomotor abnormalities. Diagnosis is clinical (second consensus criteria for MSA) with supportive features including MRI findings: putaminal atrophy, putaminal rim sign (hyperintense rim on T2-weighted MRI in MSA-P), hot cross bun sign (cruciform hyperintensity in the pons on T2-weighted MRI in MSA-C), and middle cerebellar peduncle atrophy. The putaminal rim sign on susceptibility-weighted imaging is characteristic. There is no disease-modifying treatment; management is symptomatic and supportive, including treatment of orthostatic hypotension (fludrocortisone, midodrine, droxidopa), physical therapy, speech therapy, and in advanced cases, CPAP for stridor and feeding tube placement for dysphagia. Prognosis is poor, with median survival of 6-10 years from diagnosis.

\medterm{Myasthenia Gravis}-- An autoimmune disorder caused by antibodies against the nicotinic acetylcholine receptor (AChR, 85\%) or MuSK at the neuromuscular junction, leading to fluctuating muscle weakness that worsens with use (fatigability). Presents with ptosis, diplopia, dysphagia, dysarthria, and proximal limb weakness. Ocular MG: limited to eye muscles. Generalized MG: bulbar and limb involvement. Myasthenic crisis: respiratory failure requiring ventilation. Diagnosis: AChR antibodies, repetitive nerve stimulation (decrement), single-fiber EMG, edrophonium test. Treatment: acetylcholinesterase inhibitors (pyridostigmine), corticosteroids, immunosuppressants (azathioprine, mycophenolate), thymectomy for thymoma, IVIG/plasmapheresis for crisis.

\medterm{Myelopathy}-- A general term for any neurological deficit resulting from spinal cord pathology, encompassing a wide range of causes and clinical presentations. The clinical features depend on the level and completeness of the spinal cord lesion and may include motor weakness (upper motor neuron pattern below the lesion, with spasticity, hyperreflexia, and extensor plantar responses), sensory loss (with a sensory level corresponding to the lesion), bladder and bowel dysfunction, and gait disturbance, spastic gait, sensory ataxia, or a combination. The sensory level is a key finding in myelopathy—the dermatomal level at which sensation changes from normal to abnormal below the level of the lesion. This helps localise the spinal cord lesion. Diagnosis is by MRI of the spinal cord, which demonstrates the underlying pathology (compression, inflammation, infarction, tumor, or demyelination). Management depends on the cause and may include surgical decompression for extramedullary compression (disc, tumor, abscess, haematoma), corticosteroids for inflammatory or demyelinating myelopathy, immunosuppression for autoimmune causes, and rehabilitation for functional recovery. Prognosis depends on the etiology, severity of the initial deficit, and timeliness of treatment.

\medterm{Myotonic Dystrophy}-- The most common adult-onset muscular dystrophy, an autosomal dominant multisystem disorder caused by CTG trinucleotide repeat expansions in the DMPK gene (type 1, DM1) or CCTG tetranucleotide repeat expansions in the CNBP gene (type 2, DM2) on chromosome 19q13.3. DM1 is more common and more severe than DM2. The pathophysiology involves a toxic gain-of-function mechanism, where the expanded repeat RNA forms hairpin structures that sequester RNA-binding proteins (particularly muscleblind-like proteins MBNL1 and MBNL2), leading to dysregulation of alternative splicing of many downstream genes (spliceopathy). The expanded repeats can also be transcribed into toxic proteins through repeat-associated non-ATG (RAN) translation. Clinical features of DM1 include myotonia (delayed relaxation after muscle contraction, demonstrable by grip myotonia or percussion myotonia of the thenar eminence and tongue), progressive distal weakness (worse than proximal), facial weakness with a characteristic hatchet face and ptosis, frontal balding, cataracts (posterior subcapsular iridescent cataracts), cardiac conduction abnormalities (first-degree heart block, bundle branch block, complete heart block—a leading cause of death), insulin resistance, and excessive daytime somnolence. DM2 (proximal myotonic myopathy) is milder, with predominantly proximal weakness, less severe myotonia, fewer cardiac complications, and later onset. Myotonic dystrophy exhibits anticipation (earlier and more severe disease in successive generations) due to expansion of the CTG repeat during transmission, particularly through maternal inheritance. Diagnosis is clinical, confirmed by genetic testing showing expanded CTG repeats in DMPK (DM1) or CCTG repeats in CNBP (DM2). There is no disease-modifying treatment. Management is symptomatic and includes mexiletine for myotonia, cardiac monitoring with pacemaker insertion for conduction abnormalities, treatment of cataracts and endocrine dysfunction, and respiratory support for sleep-disordered breathing. Avoid depolarising muscle relaxants (suxamethonium) due to risk of malignant hyperthermia. Prognosis in DM1 is variable; patients with severe congenital DM1 have poor outcomes, while milder adult-onset forms allow near-normal lifespan with appropriate surveillance.

\medterm{Neglect Syndrome}-- Hemispatial neglect (unilateral neglect, hemineglect) is a neurological condition characterized by failure to attend to, respond to, or orient toward stimuli on the side of space contralateral to the brain lesion, in the absence of primary sensory or motor deficits sufficient to explain the neglect. It is most commonly caused by damage to the right hemisphere, particularly the right parietal lobe (especially the inferior parietal lobule and temporoparietal junction), right thalamus, or right basal ganglia. The patient fails to respond to stimuli on the contralateral (left) side, which may include visual, auditory, tactile, or even olfactory stimuli. Patients may fail to eat food on the left side of their plate, bump into objects on the left, shave only the right side of their face, and may even deny ownership of their left limbs (somatoparaphrenia). Neglect is more severe and persistent for left-sided stimuli after right hemisphere damage due to the right hemisphere's dominant role in spatial attention (both hemispheres attend to the right, but only the right attends to the left). The contralesional side is the left side. Diagnosis includes bedside tests such as line bisection (patient marks the middle of a horizontal line, bisecting it further to the right of center), cancellation tasks (patient fails to cancel targets on the left side of the page), and drawing tasks (patient omits left-sided details when drawing a clock or a daisy). The Behavioural Inattention Test (BIT) provides a standardised assessment. Management includes visual scanning training, limb activation therapy, prism adaptation, and dopamine agonists. The condition has a significant negative impact on functional recovery and rehabilitation outcomes, with less severe forms resolving over weeks to months.

\medterm{Neurocysticercosis}-- The most common parasitic infection of the central nervous system worldwide and is caused by the larval stage (cysticercus) of the pork tapeworm Taenia solium. It is endemic in Latin America, Asia, and Africa, and is an important cause of epilepsy in endemic regions. The larvae are ingested through fecal-oral transmission (ingestion of food or water contaminated with tapeworm eggs from an infected individual), and the oncospheres cross the intestinal wall, disseminate hematogenously to the brain, where they develop into cysticerci within the brain parenchyma, subarachnoid space, ventricles, or spinal cord. Clinical presentation depends on the location and number of cysts and the host immune response. Parenchymal neurocysticercosis commonly presents with seizures (the most common manifestation in endemic regions, accounting for up to 30 percent of epilepsy in endemic areas), headache, and focal neurological deficits. Ventricular cysts can cause obstructive hydrocephalus, and subarachnoid cysts can cause arachnoiditis and communicating hydrocephalus. Diagnosis is by neuroimaging: CT shows calcified dead cysts (small hyperdense nodules) and viable cysts (hypodense lesions with a scolex visible as a hyperdense dot within the cyst). MRI is more sensitive, showing the cyst fluid intensity similar to CSF and the scolex as a mural nodule on T1-weighted imaging. Serological testing (enzyme-linked immunoelectrotransfer blot) in serum and CSF supports the diagnosis. Management includes antiparasitic drugs (albendazole, praziquantel) for active parenchymal cysts, corticosteroids to control inflammation and cerebral edema during antiparasitic treatment, anticonvulsants for seizure control, and surgical management of hydrocephalus (CSF shunting) and intraventricular or spinal cysts. Prognosis is generally good for parenchymal disease with treatment; ventricular and subarachnoid disease have worse outcomes. Prevention focuses on improved hygiene, sanitation, and meat inspection.

\medterm{Neuromyelitis Optica}-- Neuromyelitis optica spectrum disorder (NMOSD), formerly known as Devic disease, is an autoimmune, inflammatory demyelinating disorder of the central nervous system that preferentially targets the optic nerves and spinal cord. It is caused by pathogenic aquaporin-4 (AQP4) immunoglobulin G autoantibodies in approximately seventy-five to eighty percent of cases, which bind to AQP4 water channels on astrocyte foot processes, activating complement-mediated astrocyte damage and secondary demyelination. Clinical presentation of NMOSD includes optic neuritis (acute or subacute vision loss, often bilateral or sequential, and more severe than in multiple sclerosis, with poor recovery), longitudinally extensive transverse myelitis (LETM, defined as spinal cord lesion extending three or more vertebral segments on MRI, causing paraplegia, sensory loss, and bladder dysfunction), area postrema syndrome (intractable nausea, vomiting, and hiccoughs due to involvement of the medullary area postrema, which is a highly characteristic presentation), and acute brainstem syndrome. The distinction from multiple sclerosis is critical because treatments differ: MS disease-modifying therapies (interferon-beta, fingolimod, natalizumab) are ineffective or may exacerbate NMOSD. Diagnosis is based on clinical criteria (IPND 2015 criteria) requiring at least one core clinical characteristic plus positive AQP4-IgG serology. In seronegative cases, myelin oligodendrocyte glycoprotein (MOG) antibody-associated disease should be considered. MRI findings include longitudinally extensive optic nerve involvement, LETM with central cord involvement, and periaqueductal and area postrema brain lesions. Acute treatment includes high-dose intravenous corticosteroids as first-line, with plasma exchange for steroid-resistant cases. Long-term immunosuppression to prevent relapses includes mycophenolate mofetil, azathioprine, rituximab, eculizumab (a complement C5 inhibitor), inebilizumab (anti-CD19 monoclonal antibody), and satralizumab (anti-IL-6 receptor monoclonal antibody). Prognosis is relapse-dependent, with most patients accumulating disability from recurrent attacks. Attacks are typically more severe than in MS, but progression is attack-related rather than gradual.

\medterm{Neuropathic Pain}-- Chronic pain caused by damage or disease affecting the somatosensory nervous system, including the peripheral nerves, dorsal root ganglia, spinal cord, or brain. It is characteristically described as burning, shooting, electric shock-like, tingling, or stabbing, and may be accompanied by allodynia (pain from a normally non-painful stimulus) and hyperalgesia (exaggerated pain from a normally painful stimulus). Common causes include diabetic peripheral neuropathy, postherpetic neuralgia (the most common cause), painful diabetic peripheral neuropathy, trigeminal neuralgia, spinal cord injury pain, post-stroke central pain, chemotherapy-induced peripheral neuropathy, and HIV-associated neuropathy. Diagnosis is clinical, based on the characteristic pain descriptors and neurological examination findings. Screening tools such as the Douleur Neuropathique 4 (DN4) questionnaire and the Leeds Assessment of Neuropathic Symptoms and Signs (LANSS) scale help differentiate neuropathic from nociceptive pain. Nerve conduction studies and skin biopsy (quantifying intraepidermal nerve fibre density) may confirm small fibre neuropathy. Management follows a stepwise approach: first-line treatments include gabapentinoids (gabapentin, pregabalin), tricyclic antidepressants (amitriptyline, nortriptyline), and serotonin-noradrenaline reuptake inhibitors (duloxetine, venlafaxine). Topical treatments include lidocaine 5 percent patches and capsaicin 8 percent patches for localised neuropathic pain. Second-line options include tramadol and other opioid analgesics, which should be used cautiously due to the risk of dependence. Non-pharmacological treatments include transcutaneous electrical nerve stimulation (TENS), cognitive behavioural therapy, and spinal cord stimulation for refractory cases. Management of the underlying condition (e.g., glycaemic control in diabetes, antiviral therapy for herpes zoster) is also important.

\medterm{Neuropathy (Peripheral Neuropathy)}-- Damage to peripheral nerves causing sensory, motor, or autonomic dysfunction. Common causes: diabetes mellitus, alcohol use disorder, vitamin B12 deficiency, chemotherapy, autoimmune diseases, and hereditary conditions (Charcot-Marie-Tooth disease). Presents with numbness, tingling, burning pain, weakness, and gait instability. Treatment addresses the underlying cause, with symptomatic management including gabapentin, duloxetine, and tricyclic antidepressants.

\medterm{NIH Stroke Scale}-- The National Institutes of Health Stroke Scale (NIHSS) is a standardized, validated clinical assessment tool used to quantify the severity of acute ischemic stroke. It consists of fifteen items assessing level of consciousness, gaze, visual fields, facial palsy, motor strength in the arms and legs, limb ataxia, sensory function, language, dysarthria, and extinction/inattention. Each item is scored with a specific number of points, and the total score ranges from zero (no stroke symptoms) to forty-two (most severe stroke). Each item is assessed by a certified examiner, and the scale can be completed in approximately 5-8 minutes. The NIHSS is used to assess baseline stroke severity, guide treatment decisions (e.g., thrombolysis is generally recommended for NIHSS 4-25, though patients with severe strokes may still benefit), monitor neurological changes during the acute phase, and predict outcomes. Higher scores predict poorer outcomes at 90 days. The NIHSS has well-established reliability and validity but has limitations, including poor sensitivity for posterior circulation stroke and right hemisphere stroke, and language barriers for aphasic patients. It is also less accurate for very mild (NIHSS 0-3) or very severe strokes.

\medterm{Non communicating Hydrocephalus}-- Non-communicating (obstructive) hydrocephalus is a type of hydrocephalus in which there is obstruction to the flow of cerebrospinal fluid (CSF) within the ventricular system, preventing CSF from reaching the arachnoid granulations for absorption. The obstruction can occur at any level of the ventricular system, causing dilation of the ventricles proximal to the obstruction. Common causes include tumors (particularly those obstructing the fourth ventricle or the aqueduct of Sylvius, such as posterior fossa tumors such as medulloblastoma, ependymoma, and brainstem glioma), aqueductal stenosis (congenital or acquired), intraventricular hemorrhage (particularly in premature infants causing post-hemorrhagic hydrocephalus), and intraventricular cysts or masses. Clinical presentation depends on the severity and rate of progression of the obstruction. Acute obstruction presents with severe headache, nausea and vomiting, papilloedema, and rapidly decreasing level of consciousness (life-threatening). Chronic obstruction presents with headache, developmental delay (in children), macrocephaly in infants, gait disturbance, cognitive decline, and the sunsetting eye sign (upward gaze palsy due to pressure on the tectal plate from a dilated third ventricle). Diagnosis is by neuroimaging (CT or MRI) showing dilation of the ventricles proximal to the obstruction (lateral and third ventricles above the aqueduct, or fourth ventricle below the aqueduct if the obstruction is at the foramen of Luschka and Magendie), with normal or small ventricles distal to the obstruction. Aqueductal flow studies on MRI can confirm obstruction. Management includes urgent CSF diversion via endoscopic third ventriculostomy (ETV) for aqueductal stenosis and selected obstructive causes, or ventriculoperitoneal (VP) shunt placement, with removal of the obstructing lesion if possible (e.g., tumor resection).

\medterm{Normal Pressure Hydrocephalus}-- A potentially reversible cause of dementia characterized by the classic triad of gait disturbance, urinary incontinence, and cognitive impairment, with enlarged cerebral ventricles on imaging despite normal cerebrospinal fluid pressure on lumbar puncture. NPH is classified as idiopathic (primary) or secondary to conditions such as subarachnoid hemorrhage, meningitis, or traumatic brain injury. The pathophysiology is thought to involve impaired CSF absorption at the arachnoid granulations due to reduced CSF conductance, leading to increased ventricular pressure that is not reflected by a single normal lumbar puncture opening pressure but may be detected by prolonged intracranial pressure monitoring. The classic clinical triad (Wet, Wacky, Wobbly) includes gait apraxia (the earliest and most prominent feature—difficulty initiating walking with a magnetic, shuffling gait described as feet feeling stuck to the floor), cognitive impairment (subcortical dementia with psychomotor slowing, inattention, and executive dysfunction rather than the prominent memory loss of Alzheimer disease), and urinary incontinence (urinary urgency and frequency progressing to complete incontinence, typically occurring last in the triad). Diagnosis is clinical with supportive imaging: MRI brain shows enlarged ventricles out of proportion to cortical atrophy (ventriculomegaly with Evan index greater than 0.3), periventricular white matter changes, and tight high-convexity sulci (disproportionately enlarged subarachnoid space hydrocephalus). Additional diagnostic tests include a positive response to CSF drainage via lumbar puncture (Tap test) or extended lumbar CSF drainage (ELD), where improvement in gait is the most predictive finding. CSF infusion studies may show reduced CSF conductance. Management is CSF diversion via a ventriculoperitoneal (VP) shunt. Response rates are variable, with gait improvement occurring in 70-80 percent of appropriately selected patients, cognitive improvement in 50-70 percent, and urinary improvement in 50-60 percent. The best outcomes are seen in patients with dominant gait disturbance, known etiology (secondary NPH), and a positive response to CSF drainage.

\medterm{Occipital Neuralgia}-- A paroxysmal, severe, stabbing or electric shock-like pain in the distribution of the greater, lesser, or third occipital nerves (posterior scalp, from the nuchal line to the vertex). The pain is unilateral or bilateral, often described as 'ice pick' or 'stabbing', with tenderness over the occipital nerve and paraesthesias in the affected dermatome. Attacks last seconds to minutes but may be followed by a dull ache. Causes: entrapment or irritation of the occipital nerves (trauma, whiplash, cervical spine disease—C1--C3 facet arthropathy, tumors, vascular compression). Differential diagnosis: cervicogenic headache, migraine (especially occipital migraine), giant cell arteritis (must be excluded—check ESR, CRP, temporal artery biopsy), and tension-type headache. Diagnosis is clinical; occipital nerve block (injection of lidocaine and corticosteroid at the midpoint between the occipital protuberance and mastoid) is both diagnostic and therapeutic. Treatment: nerve blocks, NSAIDs, muscle relaxants, neuropathic pain agents (gabapentin, amitriptyline), and in refractory cases, pulsed radiofrequency ablation or occipital nerve stimulation.

\medterm{Oligoclonal Bands}-- Immunoglobulin G (IgG) proteins produced by a limited number of B-cell clones in the cerebrospinal fluid (CSF), detected by isoelectric focusing of CSF and serum samples. Oligoclonal bands represent intrathecal IgG synthesis and are present in approximately ninety to ninety-five percent of patients with multiple sclerosis, making them one of the most sensitive laboratory markers for the disease. The detection of OCBs in CSF but not in serum (type 2 pattern) indicates intrathecal IgG synthesis and supports the diagnosis of multiple sclerosis (MS). The presence of identical OCBs in both CSF and serum (type 3 pattern) indicates systemic IgG production with passive transfer into the CSF, which is less specific for MS. OCBs are detected in approximately 90-95 percent of MS patients, making them the most sensitive laboratory marker for the disease. They are also present in other inflammatory and infectious CNS conditions such as neuromyelitis optica, neurosarcoidosis, neurosyphilis, and neuroborreliosis, but with lower frequency. The absence of OCBs in a patient with typical MS clinical and MRI findings is associated with a milder disease course. OCB typing by isoelectric focusing of paired CSF and serum samples is considered the gold standard method.

\medterm{Pachymeningitis}-- Inflammation of the dura mater (pachymeninx), the outermost layer of the meninges. Pachymeningitis is less common than leptomeningitis (inflammation of pia and arachnoid mater). Causes: infection (tuberculosis—most common in developing countries, syphilis, fungi—Cryptococcus, bacteria—direct spread from sinusitis/otitis/osteomyelitis), autoimmune (IgG4-related disease—hypertrophic pachymeningitis, granulomatosis with polyangiitis, sarcoidosis, rheumatoid arthritis), and neoplastic (meningeal carcinomatosis, lymphoma). Presentation: chronic headache (most common), cranial nerve palsies (II, III, VI, VII, VIII), radicular pain, and myelopathy (spinal pachymeningitis). Hypertrophic pachymeningitis: focal or diffuse dural thickening on MRI with contrast enhancement (pachymeningeal enhancement). CSF: lymphocytic pleocytosis, elevated protein. Diagnosis: MRI brain with contrast, CSF analysis, and biopsy if autoimmune or neoplastic cause suspected. Treatment: corticosteroids for autoimmune causes, antitubercular therapy for TB, antibiotics for bacterial infection. Prognosis depends on the underlying etiology.

\medterm{Paranoia}-- A thought process characterized by excessive or irrational suspiciousness and distrust of others, often manifesting as delusions of persecution, reference, or grandeur. Paranoia occurs across a spectrum: paranoid personality disorder (pervasive, lifelong pattern of distrust without the severity of delusions), delusional disorder (persecutory type—non-bizarre delusions in the absence of other psychotic symptoms), paranoid schizophrenia (prominent persecutory delusions with hallucinations and disorganised thinking), and secondary paranoia (due to dementia, substance use—cocaine, amphetamines, alcohol withdrawal, traumatic brain injury, or brain tumors). Paranoia may also occur in mood disorders (psychotic depression, bipolar mania). Treatment: antipsychotic medication, cognitive behavioural therapy, and treatment of underlying cause.

\medterm{Paresthesia}-- An abnormal sensation in the skin, described as tingling, prickling, numbness, burning, or a pins-and-needles sensation, caused by dysfunction of the sensory nerves. It can be caused by any condition that affects the peripheral nerves, nerve roots, or central sensory pathways. Common causes include peripheral neuropathy (diabetes, alcohol, vitamin deficiencies), nerve compression (carpal tunnel syndrome, ulnar neuropathy, radiculopathy), multiple sclerosis, spinal cord lesions, stroke, transient ischemic attack, migraine with aura (sensory aura), anxiety and hyperventilation, and certain medications. Paresthesia is a symptom rather than a diagnosis and may be transient (e.g., pressure on a nerve causing the limb to fall asleep) or chronic (e.g., diabetic peripheral neuropathy). The distribution of the paresthesia is important for localising the lesion: peripheral nerve distribution suggests a mononeuropathy, stocking-and-glove distribution suggests a length-dependent polyneuropathy, dermatomal distribution suggests a radiculopathy, a hemi-body distribution suggests a central (brain) lesion, and bilateral leg involvement suggests a spinal cord lesion. Diagnosis involves identifying the underlying cause through history, neurological examination, electrodiagnostic studies, and appropriate laboratory tests. Management is directed at the underlying cause, with symptomatic treatment of neuropathic pain if the paresthesia is painful.

\medterm{Parinaud Syndrome}-- A clinical condition caused by lesions of the dorsal midbrain, particularly the pretectal area and superior colliculi, resulting in a characteristic triad of vertical gaze palsy (particularly impaired upward gaze), convergence-retraction nystagmus (involuntary retraction of the eyeballs into the orbits with attempted upward gaze), and light-near dissociation (pupils that do not react to light but do react to accommodation). Other features include Collier sign (eyelid retraction, also known as the lid-lag sign or setting-sun sign), convergence-retraction nystagmus, and papilloedema. Parinaud syndrome is most commonly caused by compression of the dorsal midbrain by a pineal gland tumor (pinealoma, pineal germinoma, pineocytoma, pineoblastoma), which accounts for the classic presentation in young males. Other causes include midbrain infarction, hemorrhage, multiple sclerosis, arteriovenous malformations, neurocysticercosis, and hydrocephalus (which can cause dilation of the third ventricle with upward pressure on the tectum). Diagnosis is by MRI brain with particular attention to the pineal region and dorsal midbrain, including CSF cytology and tumor markers (alpha-fetoprotein, beta-human chorionic gonadotropin) for suspected germ cell tumors. Management is directed at the underlying cause: surgical resection of pineal tumors, radiotherapy or chemotherapy for germinomas, ventriculoperitoneal shunting for hydrocephalus, and treatment of vascular or demyelinating causes. Prognosis depends on the underlying etiology and response to treatment.

\medterm{Parkinson Disease}-- A progressive neurodegenerative disorder characterized by the loss of dopaminergic neurons in the substantia nigra pars compacta of the midbrain, leading to reduced dopamine levels in the striatum and the cardinal motor features of bradykinesia, resting tremor, rigidity, and postural instability. The cardinal motor features include bradykinesia (slowness of movement, the hallmark feature, with progressive reduction in amplitude and speed of repetitive movements), resting tremor (a pill-rolling tremor of the hand at 4-6 Hz that is present at rest and diminishes with voluntary movement, typically asymmetrical at onset), rigidity (lead-pipe or cogwheel resistance to passive movement), and postural instability (impaired balance and tendency to fall, typically occurring later in the disease course). Non-motor features are increasingly recognised and include hyposmia (reduced sense of smell, often a prodromal feature that may precede motor symptoms by years), REM sleep behaviour disorder (acting out dreams), constipation, depression (the most common neuropsychiatric symptom), anxiety, apathy, cognitive impairment (executive dysfunction and bradyphrenia, progressing to dementia in up to 80 percent of patients with long-standing disease), autonomic dysfunction (orthostatic hypotension, urinary frequency, erectile dysfunction), and pain. Diagnosis is clinical using the UK Parkinson Disease Society Brain Bank criteria, which require bradykinesia plus at least one of rigidity, resting tremor, or postural instability. Dopamine transporter (DaT) SPECT scanning can confirm presynaptic dopaminergic denervation when the diagnosis is uncertain. The gold standard treatment is levodopa combined with a peripheral decarboxylase inhibitor (carbidopa or benserazide), which is the most effective symptomatic therapy, but long-term use is complicated by motor fluctuations (wearing-off phenomena, on-off fluctuations) and dyskinesias (involuntary choreiform movements). Other medications include dopamine agonists (pramipexole, ropinirole, rotigotine), monoamine oxidase B inhibitors (selegiline, rasagiline), catechol-O-methyltransferase inhibitors (entacapone, opicapone), anticholinergics (for tremor in younger patients), and amantadine (for dyskinesias). Deep brain stimulation of the subthalamic nucleus or globus pallidus internus is effective for advanced disease with motor fluctuations and dyskinesias. Multidisciplinary management including physiotherapy, occupational therapy, speech therapy, and specialist nursing is essential.

\medterm{Parkinsonism}-- A clinical syndrome characterized by bradykinesia, rigidity, tremor (resting, pill-rolling), and postural instability. Idiopathic Parkinson disease: most common cause, responds well to levodopa. Atypical parkinsonian disorders: multiple system atrophy (Parkinsonism plus cerebellar/autonomic), progressive supranuclear palsy (vertical gaze palsy, early falls), corticobasal degeneration (asymmetric apraxia, alien limb), dementia with Lewy bodies (Parkinsonism plus early dementia). Drug-induced parkinsonism: antipsychotics, resolves with drug withdrawal. Vascular parkinsonism: lower body predominance, poor levodopa response.

\medterm{Petit Mal}-- See Absence Seizure.

\medterm{Pituitary Adenoma}-- Pituitary adenomas are benign neoplasms arising from the adenohypophysis (anterior pituitary gland), accounting for approximately ten to fifteen percent of all intracranial neoplasms. They are classified by size as microadenomas (less than 10 millimeters) or macroadenomas (10 millimeters or greater), and by functional status as functioning (hormone-secreting) or non-functioning. Functioning adenomas include prolactinomas (most common functioning adenoma, causing galactorrhea, amenorrhea, and infertility. Non-functioning adenomas present with mass effect: visual field deficits (bitemporal hemianopia from compression of the optic chiasm), headache, hypopituitarism (from compression of the normal pituitary gland), and cranial nerve palsies from cavernous sinus invasion. Prolactinomas are treated medically with dopamine agonists (cabergoline is first-line, bromocriptine is second-line), which are highly effective in reducing prolactin levels and tumor size. Surgery (trans-sphenoidal adenomectomy) is indicated for non-functioning adenomas causing mass effect, functioning adenomas not responsive to medical therapy (except prolactinomas), and apoplexy (hemorrhagic infarction of a pituitary adenoma causing acute headache, visual loss, ophthalmoplegia, and altered consciousness—a pituitary emergency requiring urgent corticosteroid replacement and surgical decompression). Radiotherapy or radiosurgery (Gamma Knife) is used for residual, recurrent, or aggressive adenomas. Endocrine assessment and replacement of deficient hormones are essential before and after treatment. Prognosis is excellent for microadenomas, and macroadenomas have good outcomes with surgical resection, though long-term endocrine follow-up is required.

\medterm{Polymicrogyria}-- A brain malformation characterized by an excessive number of abnormally small, fused gyri, resulting in an irregular, bumpy cortical surface with shallow sulci. Caused by abnormal late neuronal migration or post-migrational cortical organization. May be focal or diffuse, unilateral or bilateral, and symmetric or asymmetric. Presents with variable severity depending on extent: seizures (often focal epilepsy), developmental delay, motor deficits, and language impairment. Bilateral perisylvian polymicrogyria presents with pseudobulbar palsy, dysarthria, and cognitive impairment. Diagnosis by MRI (thin-section T1-weighted and T2-weighted imaging). No specific treatment; management is symptomatic.

\medterm{Primary Lateral Sclerosis}-- A rare neurodegenerative disorder confined to the upper motor neurons (corticospinal tract) of the motor cortex and brainstem. Presents with slowly progressive spasticity, hyperreflexia, Babinski sign, and pseudobulbar affect (emotional lability). Unlike ALS, there is no significant lower motor neuron involvement (no muscle atrophy or fasciculations) and survival is longer (10-15 years vs. 2-5 years for ALS). Diagnosis: clinical, EMG/NCS to exclude ALS mimics. No disease-modifying therapy; symptomatic management with baclofen, tizanidine, and speech therapy.

\medterm{Progressive Supranuclear Palsy}-- A neurodegenerative tauopathy characterized by the accumulation of tau protein in the basal ganglia, brainstem, and cerebral cortex, belonging to the family of Parkinson plus syndromes. The cardinal clinical features include progressive vertical supranuclear gaze palsy (particularly impairment of voluntary downward gaze, which is the hallmark feature), early postural instability with recurrent unexplained falls (typically backwards), axial rigidity greater than in the limbs (axial rigidity), bradykinesia, and early cognitive decline (including frontal lobe dysfunction with apathy, disinhibition, and executive impairment, which helps distinguish PSP from Parkinson disease). Other common features include dysphagia, dysarthria (hypokinetic or spastic), pseudobulbar affect (emotional lability with pathological laughing or crying), and eyelid opening apraxia (difficulty opening the eyes voluntarily). The classic PSP phenotype is Richardson syndrome (RS), but other variants include PSP-parkinsonism (PSP-P, resembling Parkinson disease with moderate levodopa response), PSP-progressive gait freezing (PSP-PGF), and PSP-corticobasal syndrome (PSP-CBS). The pathophysiology involves abnormal aggregation of the microtubule-associated protein tau, predominantly the 4-repeat (4R) tau isoform, in neurons and glia, forming globose neurofibrillary tangles and tufted astrocytes. The H1 haplotype of the MAPT gene is the major genetic risk factor. Diagnosis is clinical (MDS-PSP criteria) with supportive MRI findings including the hummingbird sign (midbrain atrophy on sagittal T1-weighted MRI showing a flat or concave superior profile of the midbrain resembling a hummingbird) and the Mickey Mouse sign (preserved pons with atrophied midbrain on axial MRI, giving the appearance of Mickey Mouse ears). Midbrain-to-pons area ratio of less than 0.2 is highly suggestive. There is no disease-modifying treatment; management is symptomatic with levodopa for parkinsonism (limited response, usually modest and transient), botulinum toxin for blepharospasm and eyelid apraxia, speech therapy, swallowing assessment, and fall prevention strategies. Prognosis is poor, with median survival of 6-9 years from symptom onset, with falls, dysphagia, and aspiration pneumonia as common causes of death.

\medterm{Rippling Muscle Disease}-- A rare, benign myopathy characterized by mechanically-induced muscle contractions that spread across the muscle as a visible wave (rippling), typically triggered by stretch, percussion, or movement. Rippling muscle disease is caused by mutations in the CAV3 gene (encoding caveolin-3), inherited as an autosomal dominant trait with variable expressivity. Caveolin-3 is a structural protein in the sarcolemma membrane of skeletal muscle, involved in membrane stabilisation and signalling. Onset can occur at any age from childhood to adulthood. Clinical features: muscle rippling (percussion-induced or stretch-induced visible wave of muscle contraction, lasting 1--5 seconds), muscle mounding (percussion-induced localised muscle swelling, lasting up to 30 seconds), myoedema (localised dimpling on percussion), muscle stiffness/cramps after exercise, and mild proximal weakness in some patients. Serum creatine kinase is elevated (2--10 times normal). EMG shows no myotonic discharges (distinguishing from myotonic dystrophy and channelopathies). Muscle biopsy: normal or mild non-specific changes; caveolin-3 immunostaining shows reduced sarcolemmal caveolin-3. Treatment: symptomatic—avoidance of triggering stimuli, mexiletine or carbamazepine may reduce muscle stiffness. Prognosis: benign, non-progressive, normal life expectancy. Consider genetic counseling.

\medterm{Romberg Test}-- The Romberg test is a clinical neurological test used to assess the integrity of the dorsal column-medial lemniscus pathway (proprioception) and the cerebellum in maintaining balance. The test is performed by asking the patient to stand with feet together and eyes open, and then to close the eyes. A positive Romberg test is indicated by loss of balance (increased sway or falling) when the eyes are closed, with the patient being stable with eyes open. A positive result indicates impaired proprioception (dorsal column dysfunction) causing sensory ataxia, with the patient able to maintain balance when visual input compensates for the loss of proprioception. A negative Romberg test (patient remains stable with eyes closed) suggests a cerebellar cause of ataxia rather than a sensory cause. The test should be performed with the patient standing on a firm, level surface with adequate lighting and within arm's reach of a wall or examiner to prevent falls. The patient should be allowed to stand with eyes open initially, and if stable, then asked to close the eyes. The test is considered positive if the patient sways or falls within approximately 20-30 seconds of eye closure. A positive Romberg test indicates dysfunction of the dorsal columns of the spinal cord or peripheral nerves that carry proprioceptive information. Causes include tabes dorsalis (neurosyphilis), vitamin B12 deficiency (subacute combined degeneration of the cord), diabetic peripheral neuropathy, Friedreich ataxia, and other causes of sensory neuropathy.

\medterm{Schizencephaly}-- A rare brain malformation characterized by a full-thickness, gray matter-lined cleft extending from the pial surface to the lateral ventricle. Caused by early focal neuronal migration abnormalities, often associated with prenatal vascular disruption or COL4A1 mutations. Closed-lip (Type I): the cleft walls are apposed. Open-lip (Type II): there is a CSF-filled space between the cleft walls. May be unilateral or bilateral. Presents with seizures, motor deficits (hemiparesis, spasticity), developmental delay, and hydrocephalus. Diagnosis by MRI. Treatment includes seizure management, physical/occupational therapy, and hydrocephalus management if needed.

\medterm{Sciatica}-- Pain radiating along the distribution of the sciatic nerve (L4-S3), from the lower back through the buttock and down the posterior thigh to the foot. Most commonly caused by lumbar disc herniation compressing the nerve root (L4-L5 or L5-S1). Other causes: spinal stenosis, spondylolisthesis, piriformis syndrome, and epidural abscess. Presents with sharp, burning pain, paresthesias, and sometimes motor weakness (foot drop). Positive straight leg raise test. Diagnosis: clinical; MRI if red flags (cauda equina, progressive weakness, trauma, infection). Management: conservative (NSAIDs, physical therapy, activity modification); epidural steroid injections; surgical discectomy for refractory or severe cases.

\medterm{Seizure}-- A transient occurrence of signs and/or symptoms due to abnormal, excessive, or synchronous neuronal activity in the brain. Provoked seizures have an identifiable acute cause (fever, electrolyte disturbance, drug toxicity, withdrawal). Unprovoked seizures occur without an acute trigger. Classification: focal onset, generalized onset, unknown onset. Evaluation: EEG, neuroimaging, laboratory studies. Management depends on etiology and recurrence risk.

\medterm{Sleep}-- A naturally recurring physiological state characterized by altered consciousness, reduced sensory responsiveness, and distinct brain wave patterns that is essential for physical restoration, memory consolidation, and metabolic regulation. It is divided into non-rapid eye movement sleep, which includes light and deep restorative stages, and rapid eye movement sleep, which is associated with dreaming and emotional processing. Chronic sleep deprivation is linked to increased risks of obesity, cardiovascular disease, cognitive impairment, and weakened immune function, underscoring the critical role of adequate sleep in overall health.

\medterm{Sleep Studies (Polysomnography)}-- Sleep tests, most commonly polysomnography, are diagnostic studies used to evaluate sleep architecture and identify sleep disorders. Polysomnography records brain waves, eye movements, heart rate, breathing patterns, blood oxygen levels, and limb movements throughout the night in a controlled laboratory setting. Home sleep apnea tests are simpler devices that primarily measure breathing and oxygen saturation to screen for obstructive sleep apnea, while multiple sleep latency tests assess daytime sleepiness by measuring how quickly a person falls asleep in a quiet environment.

\medterm{Spasm}-- A sudden, involuntary, sustained or repetitive contraction of a muscle or group of muscles. Types: (1) Skeletal muscle spasm—muscle cramps (common, benign—nocturnal leg cramps, exercise-associated), muscle strain, tetany (hypocalcaemia—perioral paraesthesias, carpopedal spasm, Trousseau sign, Chvostek sign), tetanus (generalised rigidity, risus sardonicus, opisthotonos), dystonia (sustained muscle contractions causing twisting/repetitive movements—cervical dystonia, blepharospasm, writer's cramp), and myoclonus (brief, shock-like jerks). (2) Smooth muscle spasm—bronchospasm (asthma, COPD exacerbation), esophageal spasm (chest pain, dysphagia), ureteric colic (kidney stone—severe flank pain radiating to groin), and biliary colic (gallstone obstruction of cystic duct). Treatment: directed at the underlying cause; muscle relaxants (baclofen, tizanidine, benzodiazepines) for skeletal muscle spasm; antispasmodics (hyoscine butylbromide/Buscopan) for smooth muscle spasm; calcium channel blockers for esophageal spasm; analgesics and NSAIDs for ureteric colic.

\medterm{Spina Bifida Meningocele}-- A form of spina bifida where the meninges (but not neural tissue) herniate through a vertebral defect, forming a CSF-filled sac. The spinal cord remains in its normal position and is not involved. Presents as a visible, fluid-filled, translucent sac on the back, typically without neurologic deficits. Surgical repair involves reduction and closure. Prognosis is excellent with normal neurologic function expected.

\medterm{Spina Bifida Myelomeningocele}-- The most severe form of spina bifida, where both the meninges and neural tissue (spinal cord and nerve roots) herniate through a vertebral defect. Presents as a visible, flat or sac-like lesion on the back with neurologic deficits below the level of the lesion (paralysis, sensory loss, bowel/bladder dysfunction). Associated with hydrocephalus (in 80-90\% of cases) and Arnold-Chiari malformation type II. Diagnosis by prenatal ultrasound and elevated maternal serum alpha-fetoprotein. Treatment: surgical closure within 24-48 hours of birth, ventriculoperitoneal shunt for hydrocephalus, and lifelong multidisciplinary care.

\medterm{Spinal Shock}-- A state of transient physiological areflexia and flaccid paralysis occurring immediately after acute spinal cord injury, below the level of the lesion. It is characterized by loss of all somatic and autonomic nervous system function, including motor paralysis, sensory loss, areflexia, hypotension (from loss of sympathetic tone), bradycardia (from unopposed vagal tone), and loss of bladder and bowel function. The spinal shock phase typically begins immediately after injury and lasts days to several weeks, with gradual recovery of reflexes and autonomic function as spinal cord edema resolves and neuronal function returns. The duration of spinal shock is longer in more severe injuries. The development of hyperreflexia and spasticity marks the transition from spinal shock to the chronic phase of spinal cord injury. Management during the spinal shock phase includes aggressive hemodynamic support (maintaining mean arterial pressure >85 mmHg with IV fluids and vasopressors to prevent secondary ischemic injury to the spinal cord), immobilisation of the spine, and early surgical decompression for correctable compressive causes. The return of the bulbocavernosus reflex is the earliest sign of the end of spinal shock and is used as a clinical milestone.

\medterm{Spinocerebellar Ataxias}-- A group of autosomal dominant neurodegenerative disorders characterized by progressive cerebellar ataxia, often with additional neurological features. Over forty different SCA types have been identified, with SCA1, SCA2, SCA3 (Machado-Joseph disease, the most common worldwide), SCA6, and SCA7 being the most prevalent. Most SCAs are caused by CAG trinucleotide repeat expansions encoding polyglutamine (polyQ) tracts in the affected proteins, similar to Huntington disease, with anticipation (earlier onset in successive generations) observed in some SCA types due to repeat instability, especially with paternal transmission. Clinical features include progressive gait and limb ataxia, dysarthria, and eye movement abnormalities (nystagmus, saccadic intrusions). Additional features vary by SCA type: SCA1 includes pyramidal signs, peripheral neuropathy, and dysphagia; SCA2 includes slow saccades, peripheral neuropathy, and cognitive decline; SCA3 (Machado-Joseph disease) includes dystonia, parkinsonism, and spasticity in addition to ataxia; SCA6 is a pure cerebellar ataxia with later onset; and SCA7 includes retinal degeneration causing progressive vision loss. Diagnosis is based on clinical presentation, family history (autosomal dominant inheritance), and genetic testing identifying the specific CAG repeat expansion in the relevant gene. MRI brain shows progressive cerebellar atrophy, particularly of the vermis and cerebellar hemispheres, with brainstem atrophy in some types. There is no disease-modifying treatment; management is symptomatic and supportive, including physical and occupational therapy, speech therapy, and treatment of specific complications. Prognosis is variable but most SCAs lead to severe disability and reduced lifespan over 10-30 years.

\medterm{Stroke}-- A medical emergency caused by the interruption of blood flow to the brain, resulting in neuronal death and neurological deficits. Ischemic stroke, accounting for approximately 87 percent of cases, occurs when a thrombus or embolus blocks a cerebral artery, while hemorrhagic stroke results from rupture of a blood vessel causing bleeding into or around the brain. Rapid recognition using the FAST acronym and immediate medical intervention with thrombolytic therapy or endovascular thrombectomy can significantly reduce disability and mortality, followed by rehabilitation and secondary prevention with antiplatelet agents, anticoagulation, and risk factor modification.

\medterm{Stupor}-- A state of reduced consciousness in which the patient is unresponsive to most environmental stimuli but can be temporarily aroused by vigorous, repeated stimulation (painful stimuli, loud shouting). Stupor is distinguished from coma (no response to any stimuli, no arousal), obtundation (reduced alertness but responds to verbal stimuli), and delirium (impaired attention and cognition with fluctuating consciousness). Causes: structural brain lesions (brainstem stroke—basilar artery occlusion, thalamic stroke, intracerebral hemorrhage, subdural haematoma, brain tumor, cerebral abscess), metabolic (hypoglycaemia, hyperosmolar hyperglycaemic state, hepatic encephalopathy, uraemia, hyponatraemia, hypercalcaemia, hypothyroidism, Wernicke encephalopathy), drug intoxication (opioids, benzodiazepines, barbiturates, alcohol, antipsychotics, antidepressants, carbon monoxide poisoning), CNS infection (meningitis, encephalitis, cerebral malaria), psychiatric (catatonia, dissociative stupor—rare). Assessment: ABCs, Glasgow Coma Scale (GCS <8 indicates severe brain injury, consider airway protection), pupillary response, eye movements (doll's eye, oculocephalic reflex, cold caloric testing), motor responses, and deep tendon reflexes. Emergent CT head, blood glucose, electrolytes, toxicology screen, ABG, and lumbar puncture if indicated. Management: stabilise airway, breathing, circulation; administer naloxone if opioid overdose suspected, thiamine (before glucose) for possible Wernicke, and 50\% dextrose for hypoglycaemia. Treat underlying cause.

\medterm{Subarachnoid Hemorrhage}-- Bleeding into the subarachnoid space, the area between the arachnoid membrane and the pia mater surrounding the brain, most commonly caused by rupture of a cerebral (berry) aneurysm. Approximately eighty-five percent of non-traumatic SAH cases are caused by aneurysm rupture, with the remaining cases due to perimesencephalic (benign) SAH, arteriovenous malformations, mycotic aneurysms, and other vascular pathologies. The classic presentation is thunderclap headache—the sudden onset of the worst headache of the patient's life, reaching maximal intensity within seconds to minutes. Associated symptoms include nausea and vomiting, photophobia, nuchal rigidity (meningismus from blood in the subarachnoid space), loss of consciousness, seizures, and focal neurological deficits. The Ottawa SAH Rule and the classic presentation guide the need for urgent investigation. Diagnosis begins with a non-contrast CT head, which has 100 percent sensitivity for SAH within the first 6 hours and over 93 percent sensitivity within the first 24 hours; sensitivity declines to approximately 50 percent after one week. If CT is normal but clinical suspicion is high, a lumbar puncture is performed to detect CSF xanthochromia (yellow discolouration from bilirubin breakdown products), which is the gold standard for confirming SAH. CT angiography (CTA) or digital subtraction angiography (DSA) is used to identify the causative aneurysm. Management includes initial stabilisation in a neurosurgical critical care unit, securing the aneurysm via endovascular coiling or surgical clipping (to prevent rebleeding), management of delayed cerebral ischemia (vasospasm, which peaks at days 7-10 post-SAH, treated with nimodipine and hypertensive therapy), and management of complications including hydrocephalus (requiring external ventricular drain or VP shunt), rebleeding, and electrolyte disturbances. Prognosis is guarded, with mortality of approximately 50 percent and significant disability among survivors.

\medterm{Subdural Hematoma}-- A collection of blood between the dura mater and the arachnoid mater, caused by tearing of bridging veins that traverse the subdural space. It is classified as acute (symptoms within hours to three days), subacute (symptoms within three days to three weeks), or chronic (symptoms after three weeks), depending on the rate of bleeding and time from injury. Acute subdural hematoma is most commonly caused by severe head trauma, particularly in the elderly and alcoholics, and carries a high mortality rate (30-50 percent). Subacute and chronic subdural haematomas are more common in the elderly, alcoholics, epileptics, and patients on anticoagulation, where even minor or unrecognised head trauma can cause bleeding. Chronic subdural haematoma features a liquefied collection surrounded by an inflammatory neomembrane, which can rebleed and progressively enlarge, causing subacute or slowly progressive neurological deterioration. Clinical presentation varies: acute SDH presents with decreased consciousness, focal neurological deficit (hemiparesis), and ipsilateral pupillary dilation (from uncal herniation). Chronic SDH presents with headache, cognitive decline, gait disturbance, and personality changes, mimicking dementia or normal pressure hydrocephalus. Diagnosis is by CT head without contrast: acute SDH appears as a hyperdense, crescent-shaped collection along the convexity of the brain that does not cross the midline but can cross suture lines. Chronic SDH appears as a hypodense crescentic collection. Management: large acute SDH with mass effect requires emergency surgical evacuation via craniotomy or burr hole drainage. Chronic SDH is managed with burr hole drainage under local anesthesia or conservative observation for small asymptomatic collections. Anticoagulation reversal is critical in all cases. Prognosis: acute SDH has 30-50 percent mortality, while chronic SDH has good prognosis with treatment, though recurrence after burr hole drainage is approximately 10-20 percent.

\medterm{Syringomyelia}-- A condition characterized by the formation of a fluid-filled cyst (syrinx) within the spinal cord, which can expand and elongate over time, compressing and destroying the surrounding spinal cord tissue. The most common cause is Chiari malformation type I (see separate entry), where cerebellar tonsillar ectopia obstructs CSF flow at the foramen magnum, leading to syrinx formation. Other causes include spinal cord trauma, spinal cord tumors, arachnoiditis, and spinal dysraphism. The syrinx typically expands within the central canal of the spinal cord (dilating it) in a segmental pattern, most commonly in the cervical region, and can extend rostrally into the brainstem (syringobulbia) or caudally down the entire cord. Clinical presentation depends on the location and extent of the syrinx. Typical symptoms include a capelike distribution of dissociated sensory loss (loss of pain and temperature sensation with preserved touch and vibration sensation) over the shoulders and upper back (due to crossing spinothalamic fibers being affected first in the central cord), segmental weakness and atrophy in the arms and hands (from damage to anterior horn cells), and in advanced cases, spastic paraparesis, bladder dysfunction, and Horner syndrome (if the syrinx expands laterally). Pain is common, often described as a deep, aching, or burning sensation in the neck, shoulders, and arms. Diagnosis is by MRI of the spinal cord, which demonstrates a CSF-filled cavity within the spinal cord and any associated Chiari malformation or other pathology. Management is directed at the underlying cause: for Chiari-related syringomyelia, posterior fossa decompression (suboccipital craniectomy and cervical laminectomy with duraplasty) is the treatment of choice. For post-traumatic syringomyelia, the treatment is syrinx shunting or untethering of the spinal cord. Patients with stable, asymptomatic syringomyelia can be managed conservatively with serial MRI monitoring. The natural history is variable: some syringes remain stable while others progress insidiously over years, leading to increasing neurological deficit.

\medterm{Syrinx}-- A syrinx is a fluid-filled cavity within the spinal cord, typically containing cerebrospinal fluid or a similar fluid, that develops as a result of abnormal CSF flow dynamics within the spinal cord. Syringomyelia: a syrinx in the central spinal cord. Hydromyelia: dilation of the central canal. In practice, the terms are often used interchangeably. The most common cause of syrinx formation is Chiari malformation type I, where cerebellar tonsillar ectopia obstructs CSF flow at the foramen magnum, causing abnormal CSF pressure dynamics within the central canal of the spinal cord, leading to cyst formation. Other causes include spinal cord trauma, spinal cord tumors (especially intramedullary tumors such as ependymomas and haemangioblastomas, which can cause syrinx formation through altered CSF flow dynamics), and arachnoiditis (inflammation of the arachnoid mater due to infection, hemorrhage, or previous surgery). The syrinx may be asymptomatic if small, or may cause progressive myelopathy depending on its size and location. Symptoms include central cord syndrome with dissociated sensory loss (loss of pain and temperature with preserved touch and vibration in a cape-like distribution over the shoulders and upper back), segmental weakness and atrophy (from anterior horn cell damage), spastic paraparesis (from corticospinal tract involvement), and neuropathic pain. Diagnosis is by MRI of the spinal cord, which shows the CSF-filled cavity. Management depends on the cause and symptoms: asymptomatic syringes may be observed with serial imaging, while symptomatic syringes may require surgical treatment including posterior fossa decompression for Chiari malformations, cyst drainage or shunting, or tumor resection. Prognosis is variable; timely treatment of the underlying cause may stabilise or improve symptoms, but neurological deficits may be irreversible if the syrinx has caused significant cord damage.

\medterm{Temporal Lobe}-- One of the four major lobes of the cerebral hemisphere, located beneath the lateral fissure (Sylvian fissure) on the lateral surface of the brain, extending medially to the parahippocampal gyrus. The temporal lobe is involved in auditory processing, language comprehension, memory, and emotion. Key structures: primary auditory cortex (Heschl gyrus, Brodmann area 41/42—superior temporal gyrus), Wernicke area (posterior superior temporal gyrus—language comprehension, dominant hemisphere; lesions cause receptive aphasia/Wernicke aphasia), medial temporal lobe structures (hippocampus—memory formation, consolidation, spatial navigation; amygdala—emotion processing, fear, aggression; parahippocampal gyrus—memory and visuospatial processing). Functions: auditory perception, phonemic and semantic processing (language), declarative memory (episodic and semantic—hippocampus-dependent), facial recognition (fusiform gyrus—inferior temporal), and emotional regulation (amygdala). Pathology: temporal lobe epilepsy (most common focal epilepsy), temporal lobe tumors (glioblastoma, meningioma), herpes simplex encephalitis (hemorrhagic necrosis of the medial temporal lobes), dementia (Alzheimer disease—hippocampal atrophy is an early feature), and traumatic brain injury (temporal lobe contusions are common).

\medterm{Temporal Lobe Epilepsy}-- The most common form of focal (partial) epilepsy, arising from the medial temporal lobe structures (hippocampus, amygdala, parahippocampal gyrus). Cause: hippocampal sclerosis (most common—neuronal loss and gliosis in the CA1 and CA3 regions of the hippocampus, associated with febrile seizures in childhood), tumors (low-grade gliomas, DNET), cortical dysplasia, vascular malformations, infections (herpes simplex encephalitis, neurocysticercosis), and trauma. Seizures: focal impaired awareness seizures (complex partial seizures)—characterized by an aura (epigastric rising sensation, déjà vu, jamais vu, olfactory/gustatory hallucinations, fear), followed by behavioural arrest, automatisms (lip smacking, chewing, fumbling), and postictal confusion/amnesia. Duration: 30 seconds to 2 minutes. Seizures may secondarily generalise to bilateral tonic-clonic seizures. Interictal EEG: anterior temporal sharp waves or spikes. MRI: high-resolution coronal T2/FLAIR images show hippocampal atrophy and increased T2 signal (hippocampal sclerosis). Treatment: antiseizure medications (first-line: lamotrigine, levetiracetam, carbamazepine, oxcarbazepine). Drug-resistant epilepsy (30\%): anterior temporal lobectomy or selective amygdalohippocampectomy—seizure-free outcome in 60--80\%. Vagus nerve stimulation and responsive neurostimulation are alternatives.

\medterm{Tethered Cord Syndrome}-- A congenital or acquired condition where the spinal cord is abnormally attached (tethered) to surrounding tissues, restricting its normal upward movement within the spinal canal. Caused by a thickened filum terminale, lipomeningomyelocele, diastematomyelia, or previous myelomeningocele repair. Presents with back pain, leg weakness, sensory loss, gait abnormalities, urinary incontinence, and foot deformities. Symptoms may worsen with growth, flexion, or trauma. Diagnosis by MRI showing a low-lying conus medullaris (below L2-L3). Treatment: surgical detethering to prevent progressive neurologic deterioration.

\medterm{Thrombolysis for Stroke}-- Thrombolysis with intravenous alteplase (tPA) is the primary acute treatment for ischemic stroke when administered within the time window of symptom onset. Alteplase is a recombinant tissue plasminogen activator that converts plasminogen to plasmin, which degrades fibrin in the thrombus, restoring blood flow to the ischemic brain. The standard dose is 0.9 milligrams per kilogram (maximum 90 milligrams), with ten percent given as an intravenous bolus over one minute and the remainder infused over one hour. The therapeutic window is 0-4.5 hours from symptom onset for intravenous thrombolysis, with stricter criteria applied after 3 hours (age over 80, NIHSS >25, prior stroke plus diabetes, and anticoagulant use becoming relative contraindications). Treatment within 3 hours has the best benefit-risk profile (number needed to treat of 8 for improved outcome). Exclusion criteria include intracranial hemorrhage, recent major surgery, gastrointestinal or urinary tract bleeding, recent head trauma, uncontrolled hypertension (systolic >185 mmHg, diastolic >110 mmHg), current anticoagulation with INR >1.7, thrombocytopenia (platelets <100,000), blood glucose <2.8 or >22.2 mmol/L, and clinical presentation suggestive of subarachnoid hemorrhage even if CT is normal. The risk of symptomatic intracranial hemorrhage is approximately 6-7 percent within the first 36 hours, which limits the net benefit in patients treated later in the window. Endovascular thrombectomy with stent retrievers is indicated for large vessel occlusion (internal carotid artery, M1 segment of middle cerebral artery, basilar artery) within 6 hours of onset, and may be extended to 24 hours in selected patients based on perfusion imaging (DAWN and DEFUSE-3 trial criteria).

\medterm{Tic Douloureux}-- See Trigeminal Neuralgia.

\medterm{Transient Ischemic Attack}-- A transient ischemic attack (TIA) is a temporary episode of neurological dysfunction caused by focal brain, spinal cord, or retinal ischemia without acute infarction. Historically defined by an arbitrary time threshold of twenty-four hours, the modern tissue-based definition recognizes TIA as a brief episode of neurological symptoms caused by focal brain ischemia that resolves without evidence of acute infarction on neuroimaging (typically diffusion-weighted MRI). Most TIAs resolve within one hour, with the majority resolving within 10-60 minutes. The clinical presentation depends on the affected vascular territory and can include transient hemiparesis, hemisensory loss, aphasia, visual field defects, diplopia, vertigo, dysarthria, ataxia, and amaurosis fugax (transient monocular vision loss, the classic retinal TIA, which is a harbinger of ipsilateral carotid artery stenosis). The ABCD2 score is used to stratify the short-term risk of stroke after TIA: Age over 60 (1 point), Blood pressure >140/90 (1 point), Clinical features (unilateral weakness 2 points, speech impairment without weakness 1 point), Duration (10-59 minutes 1 point, 60+ minutes 2 points), and Diabetes (1 point). Patients with ABCD2 scores of 4 or higher or with cerebral ischemia on MRI (DWI-positive) are at higher risk and typically require urgent evaluation. The urgency of evaluation and treatment is critical because the risk of stroke is highest in the first 48 hours after a TIA, with approximately 10 percent of patients having a stroke within 90 days. Emergency management includes urgent neuroimaging (MRI brain with DWI, MRA of the head and neck), ECG and cardiac monitoring to detect atrial fibrillation, carotid Doppler ultrasound for significant carotid stenosis, and laboratory investigations for vascular risk factors. Secondary prevention includes antiplatelet therapy (aspirin plus clopidogrel for 21 days in high-risk patients, followed by single antiplatelet therapy), statin therapy (high-intensity), blood pressure control, smoking cessation, dietary modification, and carotid endarterectomy or stenting for significant ipsilateral carotid artery stenosis (70-99 percent). Anticoagulation is indicated for TIA associated with atrial fibrillation.

\medterm{Transverse Myelitis}-- An inflammatory disorder of the spinal cord causing acute or subacute motor, sensory, and autonomic dysfunction. It results from immune-mediated inflammation of the spinal cord parenchyma and can occur as a isolated condition or as part of a systemic autoimmune disease or demyelinating disorder such as multiple sclerosis or neuromyelitis optica. Clinical features include bilateral lower extremity weakness, a sensory level (a dermatome below which sensation is impaired), urinary retention (more common than incontinence in the acute phase, progressing to bladder and bowel incontinence), and a sensory level (a transverse band-like boundary below which there is loss of pain, temperature, and light touch sensation). The sensory level is typically in the mid-thoracic region, as the thoracic spinal cord is the longest and most vulnerable segment. Fever or a recent history of infection or vaccination may precede the onset. If the inflammation occurs as part of multiple sclerosis, it is typically shorter (less than one vertebral segment) and more asymmetrical, whereas NMOSD-related transverse myelitis is typically longitudinally extensive (three or more vertebral segments). Diagnosis is by MRI of the spinal cord, which shows a central T2-hyperintense intramedullary lesion with gadolinium enhancement in the acute phase. CSF analysis typically shows a lymphocytic pleocytosis (elevated white cell count) and elevated protein, with positive oligoclonal bands if associated with multiple sclerosis. Treatment of acute attacks includes high-dose intravenous corticosteroids (methylprednisolone 1 gram daily for 5 days), with plasma exchange for steroid-resistant cases. Long-term treatment depends on the underlying cause: immunosuppression for NMOSD, disease-modifying therapy for MS, and observation for monophasic idiopathic cases. Prognosis is variable: approximately one-third of patients recover fully, one-third have moderate residual disability, and one-third have severe disability.

\medterm{Trigeminal Neuralgia}-- A chronic pain condition characterized by recurrent, sudden, severe, electric shock-like pain in the distribution of the trigeminal nerve (typically V2/V3). Often triggered by light touch, chewing, or brushing teeth. Most cases are caused by vascular compression of the trigeminal nerve root. First-line treatment: carbamazepine or oxcarbazepine. Surgical options: microvascular decompression, gamma knife radiosurgery, percutaneous rhizotomy.

\medterm{Tuberous Sclerosis Complex}-- An autosomal dominant genetic disorder caused by mutations in TSC1 (hamartin) or TSC2 (tuberin) genes, leading to mTOR pathway activation. Features: cortical tubers, subependymal nodules, cardiac rhabdomyomas, angiomyolipomas, facial angiofibromas, and seizures (often infantile spasms). Intellectual disability and autism spectrum disorder are common. Treatment: mTOR inhibitors (everolimus) for angiomyolipomas and subependymal giant cell astrocytomas.

\medterm{Unconsciousness}-- See Coma.

\medterm{Vertebrobasilar Insufficiency}-- A condition of reduced blood flow through the vertebral arteries or basilar artery to the posterior circulation of the brain, including the brainstem, cerebellum, occipital lobes, and thalami. It is most commonly caused by atherosclerosis of the vertebral arteries or basilar artery, with risk factors including hypertension, diabetes, smoking, hyperlipidemia, and advanced age. The posterior circulation supplies vital structures including the brainstem (containing cranial nerve nuclei and vital autonomic centers), the cerebellum (controlling coordination and balance), the occipital lobes (vision), and the thalami (sensory and motor relay). Clinical presentation of VBI includes dizziness and vertigo (the most common symptom), diplopia, dysarthria, dysphagia, ataxia, bilateral weakness or sensory loss, drop attacks (sudden loss of postural tone without loss of consciousness), visual disturbances (blurred vision, visual field defects, cortical blindness), and tinnitus. The diagnosis of VBI is distinguished from benign paroxysmal positional vertigo (which has positional triggers), Meniere disease (which has episodic vertigo with hearing loss and tinnitus), and migraine-associated vertigo. Diagnosis is by clinical assessment of posterior circulation symptoms and signs, with imaging including MRA or CTA of the vertebral and basilar arteries to assess for stenosis, occlusion, or dissection. Ultrasound of the vertebral arteries (extracranial segment) and transcranial Doppler of the basilar artery may demonstrate flow reduction. Management includes aggressive risk factor modification, antiplatelet therapy (aspirin or clopidogrel), and revascularisation (vertebral artery angioplasty and stenting, or vertebral artery endarterectomy) for selected patients with symptomatic severe stenosis refractory to medical therapy. Anticoagulation is indicated for embolism from the heart or proximal vertebral artery. Prognosis is poorer than for anterior circulation disease, with a higher risk of disabling or fatal stroke if large vessel disease is present.

\medterm{Vestibular Neuritis}-- Acute vestibular dysfunction due to inflammation of the vestibular nerve, typically viral in origin. Presents with acute onset of vertigo, nausea, vomiting, and nystagmus lasting days. Hearing is normal (unlike labyrinthitis). Diagnosis: clinical, HINTS exam (Head Impulse, Nystagmus, Test of Skew) to differentiate from central causes. Treatment: corticosteroids if early, vestibular suppressants acutely, and vestibular rehabilitation.

\medterm{Vestibular Schwannoma}-- Vestibular schwannoma, also known as acoustic neuroma, is a benign, slow-growing tumor arising from the Schwann cells of the vestibular portion of the eighth cranial nerve (vestibulocochlear nerve), typically from the superior vestibular nerve. It is the most common tumor of the cerebellopontine angle, accounting for approximately eighty percent of tumors in this location. The most common presenting symptom is unilateral sensorineural hearing loss, which develops gradually over months to years and is often accompanied by tinnitus, vertigo or imbalance, and a sensation of fullness in the ear. Larger tumors compress the trigeminal nerve (causing facial numbness and paraesthesiae), the facial nerve (causing facial weakness or spasm), and the brainstem and cerebellum (causing ataxia, long tract signs, and hydrocephalus from compression of the fourth ventricle). The diagnosis is confirmed by MRI of the internal auditory meatus and brain with gadolinium showing an enhancing mass in the cerebellopontine angle that extends into and expands the internal auditory canal. The differential diagnosis includes meningioma, facial nerve schwannoma, epidermoid cyst, and trigeminal schwannoma. Audiometry shows unilateral sensorineural hearing loss with poor speech discrimination. Auditory brainstem response testing shows delayed latencies. Management options include observation with serial imaging for small tumors in elderly or medically unfit patients, fractionated stereotactic radiosurgery (Gamma Knife or linear accelerator) for small to medium tumors (less than 3 cm), and microsurgical resection via a translabyrinthine, retrosigmoid, or middle cranial fossa approach for large tumors causing significant mass effect or brainstem compression. The goal of radiosurgery is tumor control (prevention of growth) rather than complete elimination. Hearing preservation is possible in selected cases with small tumors and good pre-treatment hearing, but the likelihood of hearing preservation decreases with tumor size.

\medterm{Wilson Disease}-- An autosomal recessive inherited disorder of copper metabolism caused by mutations in the ATP7B gene on chromosome 13, which encodes a copper-transporting ATPase in hepatocytes. The defective ATP7B protein impairs incorporation of copper into apoceruloplasmin and biliary excretion of copper, leading to progressive copper accumulation in the liver, brain, cornea, and other organs. Clinical manifestations typically present between ages five and thirty-five and include hepatic disease (ranging from asymptomatic elevation of liver enzymes to acute hepatitis, cirrhosis, and fulminant hepatic failure), neurological disease (dysarthria, ataxia, dystonia, parkinsonism, tremor—particularly a wing-beating tremor of the arms, chorea, and cognitive impairment with executive dysfunction), psychiatric symptoms (depression, personality changes, psychosis), and Kayser-Fleischer rings (golden-brown discolouration at the limbus of the cornea due to copper deposition in Descemet membrane, visible on slit-lamp examination and pathognomonic for Wilson disease when present). Diagnosis is based on a combination of clinical features, Kayser-Fleischer rings on slit-lamp examination, low serum caeruloplasmin (less than 0.2 g/L), elevated 24-hour urinary copper excretion (greater than 100 micrograms), elevated hepatic copper content on liver biopsy, and genetic testing for ATP7B mutations. Serum free copper (non-caeruloplasmin-bound copper) is elevated. Treatment involves lifelong copper chelation therapy with D-penicillamine (first-line, acts by chelating copper and promoting urinary excretion) or trientine, zinc acetate (which blocks intestinal absorption of copper and induces hepatic metallothionein, used as maintenance therapy or first-line for presymptomatic patients), and dietary restriction of copper-rich foods (shellfish, nuts, chocolate, mushrooms, organ meats). Liver transplantation is curative for severe hepatic failure or refractory neurological disease. Prognosis is excellent with early diagnosis and lifelong treatment; without treatment, Wilson disease is progressive and fatal. Patients who stop chelation therapy are at risk of irreversible neurological deterioration and hepatic decompensation.
